% !TEX root = print.tex

\part{Approximation spaces}
\label{part:approximation}

This part develops the abstract theory: approximation schemes in a normed space, the
approximation spaces $A^s_q$ and their embeddings, Jackson's and Bernstein's inequalities with
the Nikolskii transfer theorem, the representation of $A^s_q$ as a real interpolation space of
the discrete $K$-method, block decompositions and projections, traces and intrinsic norms, and
moduli of smoothness. Its last three chapters apply the theory to schemes of functions on a
finite measure space, whose approximation spaces are Besov spaces with embeddings from a
Nikolskii inequality, and to spectral schemes in a Hilbert space. The characterisations proved
here are the abstract ones: by the error of best approximation, by blocks, by real interpolation
(the discrete $K$-method) and by moduli of smoothness.

\chapter{Approximation schemes}
\label{ch:schemes}

An approximation scheme is an increasing family of sets $A_0 \subseteq A_1 \subseteq \cdots$ in a
normed space, and $f$ belongs to the approximation space $A^s_q$ in the sense of
Pietsch~\cite{Pietsch} when its errors of best approximation $E_n(f) = \dist(f, A_n)$ decay like
$n^{-s}$ in the $\ell_q$ sense.

Throughout, $X$ and $Y$ are real normed spaces. Complex spaces, such as $L_p$ spaces of
complex-valued functions, are regarded as real ones.

\begin{definition}[Approximation scheme]
\label{def:scheme}
\lean{ApproximationScheme, ApproximationScheme.coe_zero, ApproximationScheme.ext}
\leanok
An \emph{approximation scheme} in $X$ is a sequence of subsets
$A_0 \subseteq A_1 \subseteq A_2 \subseteq \cdots$ of $X$ such that
\begin{enumerate}
  \item $A_0 = \{0\}$,
  \item $c\,a \in A_n$ for all $c \in \R$, $n \in \N$ and $a \in A_n$,
  \item $a + b \in A_{m+n}$ for all $m, n \in \N$, $a \in A_m$ and $b \in A_n$.
\end{enumerate}
\end{definition}

The elements of $A_n$ are the approximants of complexity $n$. The sets need not be linear
subspaces: in the $n$-term scheme of \Cref{def:nterm} they are unions of subspaces. A family
indexed by a degree enters with a strict inequality, $A_n = \{\text{degree} < n\}$, so that
$A_0 = \{0\}$ holds.

\begin{lemma}[Elementary properties]
\label{lem:scheme-basic}
\lean{ApproximationScheme.zero_mem, ApproximationScheme.nonempty, ApproximationScheme.neg_mem,
  ApproximationScheme.sub_mem, ApproximationScheme.sum_mem, ApproximationScheme.smul_set_eq,
  ApproximationScheme.sub_mem_pow_add_two, ApproximationScheme.mem_of_le,
  ApproximationScheme.add_mem, ApproximationScheme.add_mem_pow_succ,
  ApproximationScheme.sub_mem_pow_succ, ApproximationScheme.iUnionSubmodule,
  ApproximationScheme.coe_iUnionSubmodule}
\leanok
\uses{def:scheme}
Let $A$ be an approximation scheme in $X$. Then for all $m, n, k \in \N$:
\begin{enumerate}
  \item $0 \in A_n$; in particular $A_n$ is nonempty;
  \item $a - b \in A_{m+n}$ for all $a \in A_m$ and $b \in A_n$;
  \item $\sum_{i \in I} a_i \in A_{\sum_{i \in I} n_i}$ for every finite family with
    $a_i \in A_{n_i}$;
  \item $c\,A_n = A_n$ for every $c \in \R \setminus \{0\}$;
  \item $a - b \in A_{2^{k+2}}$ for all $a \in A_{2^{k+1}}$ and $b \in A_{2^k}$;
  \item $\bigcup_n A_n$ is a linear subspace of $X$.
\end{enumerate}
\end{lemma}

\begin{proof}
\leanok
The first item follows from $A_0 = \{0\} \subseteq A_n$, the second from $-b = (-1)\,b$, and
the third by induction over $I$. For the fourth, $c^{-1} a \in A_n$ for $a \in A_n$. The fifth
follows from the second and $2^{k+1} + 2^k \leq 2^{k+2}$, and the last from the axioms
$A_m + A_n \subseteq A_{m+n}$ and $c\,A_n \subseteq A_n$.
\end{proof}

\begin{definition}[Preimage and image of a scheme]
\label{def:comap}
\lean{ApproximationScheme.comap, ApproximationScheme.map, ApproximationScheme.ofSubmodules,
  ApproximationScheme.ofMonotoneSubmodules, ApproximationScheme.mem_ofMonotoneSubmodules_zero,
  ApproximationScheme.mem_ofMonotoneSubmodules,
  ApproximationScheme.mem_ofMonotoneSubmodules_two_pow, ApproximationScheme.mem_comap,
  ApproximationScheme.mem_map, ApproximationScheme.mem_map_of_mem,
  ApproximationScheme.mem_ofSubmodules}
\leanok
\uses{def:scheme}
Let $A$ be an approximation scheme in $X$. For an injective linear map $\iota : Y \to X$, the
sets $\iota^{-1}(A_n)$ form an approximation scheme in $Y$. For a linear map $T : X \to Y$, the
sets $T(A_n)$ form an approximation scheme in $Y$. An increasing sequence of linear subspaces
$V_0 \subseteq V_1 \subseteq \cdots$ of $X$ with $V_0 = \{0\}$ is an approximation scheme. For
every increasing sequence of linear subspaces $V_0 \subseteq V_1 \subseteq \cdots$, the sets
$A_0 = \{0\}$ and $A_n = V_{\lfloor \log_2 n \rfloor}$, $n \geq 1$, form an approximation scheme,
the \emph{dyadic scheme} of $(V_k)_k$, with $A_{2^k} = V_k$.
\end{definition}

Injectivity of $\iota$ is what gives $\iota^{-1}(\{0\}) = \{0\}$.

\section{The error of best approximation}

\begin{definition}[Error of best approximation]
\label{def:error}
\lean{ApproximationScheme.approxError, ApproximationScheme.approxError_nonneg,
  ApproximationScheme.approxError_le_norm_sub}
\leanok
\uses{def:scheme}
Let $A$ be an approximation scheme in $X$. For $n \in \N$ and $f \in X$,
\[
  E_n(f) = \dist(f, A_n) = \inf_{a \in A_n} \|f - a\|.
\]
\end{definition}

\begin{lemma}[Calculus of the error]
\label{lem:error-basic}
\lean{ApproximationScheme.approxError_zero_left, ApproximationScheme.approxError_le_norm,
  ApproximationScheme.approxError_anti, ApproximationScheme.approxError_smul,
  ApproximationScheme.approxError_add_le, ApproximationScheme.abs_approxError_sub_le,
  ApproximationScheme.approxSeq_le_approxSeq_add_enorm, ApproximationScheme.approxError_zero_right,
  ApproximationScheme.approxError_eq_zero_of_mem,
  ApproximationScheme.approxError_eq_zero_of_mem_closure,
  ApproximationScheme.approxError_anti_scheme, ApproximationScheme.le_approxError_iff,
  ApproximationScheme.exists_mem_norm_sub_lt, ApproximationScheme.exists_mem_norm_sub_le_add,
  ApproximationScheme.le_mul_approxError_of_forall_mem,
  ApproximationScheme.approxError_add_le_pow_succ, ApproximationScheme.approxError_neg,
  ApproximationScheme.approxError_le_approxError_add_norm,
  ApproximationScheme.continuous_approxError, ApproximationScheme.approxSeq_anti,
  ApproximationScheme.approxSeq_zero_left, ApproximationScheme.approxSeq_zero_right,
  ApproximationScheme.approxSeq_le_enorm, ApproximationScheme.approxSeq_ne_top,
  ApproximationScheme.approxSeq_le_of_mem, ApproximationScheme.approxSeq_eq_zero_of_mem,
  ApproximationScheme.approxSeq_add_le, ApproximationScheme.approxSeq_add_le_pow_succ,
  ApproximationScheme.approxSeq_smul, ApproximationScheme.approxSeq_neg,
  ApproximationScheme.approxSeq_anti_scheme}
\leanok
\uses{def:error, lem:scheme-basic}
Let $A$ be an approximation scheme in $X$. For all $f, g \in X$, $c \in \R$ and
$m, n \in \N$:
\begin{enumerate}
  \item $E_0(f) = \|f\|$ and $E_n(f) \leq \|f\|$;
  \item $E_n(f) \leq E_m(f)$ whenever $m \leq n$;
  \item $E_n(c f) = |c| \, E_n(f)$;
  \item $E_{m+n}(f + g) \leq E_m(f) + E_n(g)$;
  \item $|E_n(f) - E_n(g)| \leq \|f - g\|$, so $E_n$ is continuous;
  \item $E_n(f) = 0$ for $f$ in the closure of $A_n$, in particular $E_n(0) = 0$;
  \item a scheme with larger sets has smaller errors;
  \item for every $\varepsilon > 0$ there is $a \in A_n$ with
    $\|f - a\| \leq E_n(f) + \varepsilon$.
\end{enumerate}
The same statements hold for the errors as elements of $[0, \infty]$, the form used by the
approximation spaces.
\end{lemma}

\begin{proof}
\leanok
The first item holds because $A_0 = \{0\}$ and $0 \in A_n$, the second because the sets
increase, the third because $c A_n = A_n$ for $c \neq 0$. For the fourth, given
$\varepsilon > 0$ choose $a \in A_m$ and $b \in A_n$ with $\|f - a\| \leq E_m(f) +
\varepsilon$ and $\|g - b\| \leq E_n(g) + \varepsilon$; then $a + b \in A_{m+n}$ and the
triangle inequality gives the bound up to $2\varepsilon$. The fifth item is the
$1$-Lipschitz property of the distance to a set, the sixth holds because the distance to a set
equals the distance to its closure, the seventh because the infimum over a larger set is
smaller, and the last is the definition of the infimum.
\end{proof}

A best approximant, an $a \in A_n$ with $\|f - a\| = E_n(f)$, need not exist, since the sets
$A_n$ need not be closed. The proofs work with near-best approximants as in the proof of the
fourth item.

\begin{lemma}[Convergence of the errors]
\label{lem:error-tendsto}
\lean{ApproximationScheme.tendsto_approxError_zero_iff,
  ApproximationScheme.tendsto_approxError_zero_of_pow}
\leanok
\uses{lem:error-basic}
Let $A$ be an approximation scheme in $X$ and $f \in X$. Then $E_n(f) \to 0$ as
$n \to \infty$ if and only if $f$ lies in the closure of $\bigcup_n A_n$, and this holds as
soon as $E_{2^k}(f) \to 0$ as $k \to \infty$.
\end{lemma}

\begin{proof}
\leanok
If $\|f - a\| < \varepsilon$ for some $a \in A_m$, then $E_n(f) < \varepsilon$ for all
$n \geq m$ since the errors decrease. Conversely $E_n(f) < \varepsilon$ gives such an $a$.
The errors decrease, so convergence along $2^k$ gives convergence.
\end{proof}

\begin{lemma}[Bounds by a modulus extend from a dense set]
\label{lem:error-dense}
\lean{ApproximationScheme.approxSeq_le_of_dense}
\leanok
\uses{lem:error-basic}
Let $A$ be an approximation scheme in $X$, $n \in \N$, $D \subseteq X$ dense,
$C, C_0 \in [0, \infty)$, and let $\omega : X \to [0, \infty]$ satisfy
$\omega(f + g) \leq \omega(f) + \omega(g)$ and $\omega(f) \leq C_0 \|f\|$ for all
$f, g \in X$. If $E_n(g) \leq C \omega(g)$ for every $g \in D$, then $E_n(f) \leq C \omega(f)$
for every $f \in X$.
\end{lemma}

\begin{proof}
\leanok
For $g \in D$, $E_n(f) \leq E_n(g) + \|f - g\|$ and $\omega(g) \leq \omega(f) + C_0 \|f - g\|$,
so $E_n(f) \leq C \omega(f) + (C C_0 + 1) \|f - g\|$, and $\|f - g\|$ is arbitrarily small.
\end{proof}

\begin{lemma}[Functoriality of the error]
\label{lem:error-map}
\lean{ApproximationScheme.approxError_map_le}
\leanok
\uses{def:error}
Let $A$ and $B$ be approximation schemes in $X$ and $Y$, and let $T : X \to Y$ be a bounded
linear map with $T(A_n) \subseteq B_n$ for all $n$. Then for all $n \in \N$ and $f \in X$,
\[
  E^B_n(T f) \leq \|T\| \, E^A_n(f).
\]
\end{lemma}

\begin{proof}
\leanok
For $a \in A_n$, $E^B_n(Tf) \leq \|Tf - Ta\| \leq \|T\| \, \|f - a\|$; take the infimum
over $a$.
\end{proof}

\section{Two schemes}

\begin{definition}[Linear approximation]
\label{def:ofspan}
\lean{ApproximationScheme.ofSpan, ApproximationScheme.ofSpan_apply,
  ApproximationScheme.finiteDimensional_span_ofSpan}
\leanok
\uses{def:scheme}
For a sequence $(e_i)_{i \in \N}$ in $X$, the sets $A_n = \spn\{e_i : i < n\}$ form an
approximation scheme.
\end{definition}

\begin{lemma}
\label{lem:ofspan}
\lean{ApproximationScheme.isClosed_ofSpan, ApproximationScheme.iUnion_ofSpan,
  ApproximationScheme.tendsto_approxError_ofSpan_iff}
\leanok
\uses{def:ofspan, lem:error-tendsto}
For the scheme of \Cref{def:ofspan}, every $A_n$ is closed, $\bigcup_n A_n = \spn\{e_i : i
\in \N\}$, and $E_n(f) \to 0$ if and only if $f$ lies in the closed span of $(e_i)$.
\end{lemma}

\begin{proof}
\leanok
The sets are finite-dimensional subspaces and therefore closed. The union of an increasing
family of spans is the span of the union. The last statement is \Cref{lem:error-tendsto}.
\end{proof}

\begin{definition}[$n$-term approximation]
\label{def:nterm}
\lean{ApproximationScheme.nTerm, ApproximationScheme.mem_nTerm,
  ApproximationScheme.mem_nTerm_of_mem}
\leanok
\uses{def:scheme}
For a set $D \subseteq X$, the \emph{dictionary}, let $A_n$ be the set of linear combinations
of at most $n$ elements of $D$. These sets form an approximation scheme.
\end{definition}

The third axiom holds because a combination of $m$ and one of $n$ elements of $D$ together
use at most $m + n$ elements.

\begin{lemma}
\label{lem:nterm-ofspan}
\lean{ApproximationScheme.ofSpan_subset_nTerm, ApproximationScheme.approxError_nTerm_le_ofSpan}
\leanok
\uses{def:ofspan, def:nterm, lem:error-basic}
For a sequence $(e_i)$ in $X$, the linear scheme of $(e_i)$ is contained in the $n$-term
scheme of the dictionary $\{e_i : i \in \N\}$, and hence the $n$-term errors are at most the
linear ones.
\end{lemma}

\begin{proof}
\leanok
An element of $\spn\{e_i : i < n\}$ uses the $n$ elements $e_0, \dots, e_{n-1}$.
\end{proof}

\chapter{Sequence spaces}
\label{ch:sequences}

The quasi-norms of this blueprint take values in $[0, \infty]$, and a function belongs to a
space when its quasi-norm is finite. Every series of elements of $[0, \infty]$ converges, so
no summability conditions are needed. Arithmetic in $[0, \infty]$ uses the conventions
$0 \cdot \infty = 0$, $0^{-1} = \infty$, $\infty^{-1} = 0$, and $x - y = 0$ for $x \leq y$.
With these, a constant $(1 - 2^{-\alpha})^{-1}$ is finite exactly for $\alpha > 0$.

\begin{definition}[$\ell_q$ norm and dyadic weights]
\label{def:lq-norm}
\lean{ENNReal.lqNorm, ENNReal.lqNorm_top, ENNReal.lqNorm_eq_tsum, ENNReal.dyadicWeight,
  ENNReal.aestronglyMeasurable_count, ENNReal.lqNorm_one, ENNReal.one_le_toReal_of_one_le,
  ENNReal.lqNorm_rpow, ENNReal.lqNorm_two_sq, ENNReal.lqNorm_zero_right, ENNReal.lqNorm_zero_left,
  ENNReal.dyadicWeight_zero_right, ENNReal.dyadicWeight_zero_left, ENNReal.dyadicWeight_ne_top,
  ENNReal.dyadicWeight_lt_top, ENNReal.dyadicWeight_pos, ENNReal.dyadicWeight_ne_zero,
  ENNReal.dyadicWeight_add, ENNReal.dyadicWeight_add_nat, ENNReal.dyadicWeight_succ,
  ENNReal.dyadicWeight_neg, ENNReal.dyadicWeight_neg_mul_self, ENNReal.dyadicWeight_mul_neg_self,
  ENNReal.dyadicWeight_rpow, ENNReal.dyadicWeight_eq_pow, ENNReal.dyadicWeight_mono,
  ENNReal.dyadicWeight_le_of_le, ENNReal.dyadicWeight_anti, ENNReal.one_le_dyadicWeight,
  ENNReal.natCast_two_pow_rpow, ENNReal.two_pow_rpow_eq_dyadicWeight,
  ENNReal.two_rpow_neg_lt_one, ENNReal.one_sub_two_rpow_neg_inv_ne_top,
  ENNReal.one_le_one_sub_two_rpow_neg_inv, ENNReal.lqNorm_dyadicWeight_neg_le,
  ENNReal.lqNorm_lt_top_iff_of_le_mul}
\leanok
For $q \in (0, \infty)$, a countable set $I$ and a family $a = (a_i)_{i \in I}$ in $[0, \infty]$,
\[
  \|a\|_{\ell_q} = \Bigl(\sum_{i \in I} a_i^q\Bigr)^{1/q},
  \qquad
  \|a\|_{\ell_\infty} = \sup_{i \in I} a_i .
\]
For $s \in \R$ and $k \in \N$, the \emph{dyadic weight} is $2^{ks}$.
\end{definition}

In Lean, $\|a\|_{\ell_q}$ is the $L_q$ norm of $a$ for the counting measure on $I$. The index set
is $\N$ for sequences, $\N^d$ for dominating mixed smoothness, and a finite set for the
coefficients of one level of atoms and wavelets.

\begin{lemma}[Properties of the $\ell_q$ norm]
\label{lem:lq-basic}
\lean{ENNReal.lqNorm_mono, ENNReal.le_lqNorm, ENNReal.lqNorm_const_mul, ENNReal.lqNorm_add_le,
  ENNReal.lqNorm_le_lqNorm_of_le, ENNReal.lqNorm_le_head_add_shift,
  ENNReal.lqNorm_le_sum_range_add_shift, ENNReal.lqNorm_le_lqNorm_one, ENNReal.lqNorm_geometric_le,
  ENNReal.lqNorm_shift_le, ENNReal.lqNorm_eq_sum, ENNReal.lqNorm_le_card_rpow_mul,
  ENNReal.lqNorm_prod_rpow, ENNReal.lqNorm_le_sum_of_eq_zero, Rd.lqNorm_ite_eq,
  ENNReal.lqNorm_ite_le_eq}
\leanok
\uses{def:lq-norm}
Let $1 \leq q \leq r \leq \infty$. For all sequences $a, b$ in $[0, \infty]$, all
$c, \rho \in [0, \infty]$ and all $n \in \N$:
\begin{enumerate}
  \item if $a_k \leq b_k$ for every $k$, then $\|a\|_{\ell_q} \leq \|b\|_{\ell_q}$;
  \item $a_n \leq \|a\|_{\ell_q}$ and $\|c\,a\|_{\ell_q} = c \, \|a\|_{\ell_q}$;
  \item $\|a + b\|_{\ell_q} \leq \|a\|_{\ell_q} + \|b\|_{\ell_q}$;
  \item $\|a\|_{\ell_r} \leq \|a\|_{\ell_q} \leq \|a\|_{\ell_1}$;
  \item $\|a\|_{\ell_q} \leq a_0 + \|(a_{k+1})_k\|_{\ell_q}$, and
    $\|a\|_{\ell_q} \leq a_0 + \cdots + a_{m-1} + \|(a_{k+m})_k\|_{\ell_q}$ for every $m \in \N$;
  \item $\|(\rho^k)_k\|_{\ell_q} \leq (1 - \rho)^{-1}$;
  \item if $a_k = 0$ for every $k$ outside a finite set $s \subseteq \N$, then
    $\|a\|_{\ell_q} \leq \sum_{k \in s} a_k$, and the sequence equal to $c$ at $n$ and to $0$
    elsewhere has $\ell_q$ norm $c$;
  \item if $q < \infty$, the tail $(0, \dots, 0, a_n, a_{n+1}, \dots)$ has the $\ell_q$ norm
    $\|(a_{n+k})_k\|_{\ell_q}$.
\end{enumerate}
\end{lemma}

\begin{proof}
\leanok
The first two items are immediate. The third is Minkowski's inequality for the finite
partial sums, followed by the supremum over them. For the fourth put $M = \|a\|_{\ell_q}$;
then $a_k^r = a_k^{r-q} a_k^q \leq M^{r-q} a_k^q$, and summation gives
$\|a\|_{\ell_r}^r \leq M^r$. The fifth is the third applied to $a_0 \delta_0$ and the rest,
and induction on $m$. The sixth is the fourth and the geometric series, the seventh is the
fourth for a sequence with finitely many nonzero terms, and the last compares the series of the
$q$-th powers term by term.
\end{proof}

\begin{lemma}[The Cauchy--Schwarz inequality for sequences]
\label{lem:lq-cauchy-schwarz-sequences}
\lean{ENNReal.lqNorm_one_mul_le, ENNReal.sum_mul_le_lqNorm_two_mul}
\leanok
\uses{def:lq-norm}
For every countable index set $I$, every finite set $s \subseteq I$ and all families
$(a_k)_k$ and $(b_k)_k$ in $[0, \infty]$ indexed by $I$,
\[
  \sum_{k \in s} a_k b_k \leq \sum_{k \in I} a_k b_k
  \leq \|(a_k)_k\|_{\ell_2} \, \|(b_k)_k\|_{\ell_2} .
\]
\end{lemma}

\begin{proof}
\leanok
The second inequality is Hölder's inequality for the counting measure on $I$ (Mathlib's
\texttt{ENNReal.\allowbreak lintegral\_\allowbreak mul\_\allowbreak le\_\allowbreak
Lp\_\allowbreak mul\_\allowbreak Lq}) with exponents $2$ and $2$; the first holds since every
term is nonnegative.
\end{proof}

\begin{lemma}[Countable Minkowski inequality and monotone convergence in $\ell_q$]
\label{lem:countable-minkowski}
\lean{ENNReal.lqNorm_sum_le, ENNReal.lqNorm_iSup, ENNReal.lqNorm_tsum_le,
  ENNReal.monotone_ite_lt, ENNReal.lqNorm_eq_iSup_lqNorm_ite}
\leanok
\uses{lem:lq-basic}
For every $1 \leq q \leq \infty$ and every family $(f_{i,k})_{i, k \in \N}$ in $[0, \infty]$,
\[
  \Bigl\| \Bigl(\sum_{i=0}^\infty f_{i,k}\Bigr)_k \Bigr\|_{\ell_q}
  \leq \sum_{i=0}^\infty \|(f_{i,k})_k\|_{\ell_q} .
\]
For every $q \in (0, \infty]$ and every sequence $a$ in $[0, \infty]$, the truncations
$(a_k \mathbf 1_{\{k < J\}})_k$ increase with $J$ and
$\|a\|_{\ell_q} = \sup_{J \in \N} \|(a_k \mathbf 1_{\{k < J\}})_k\|_{\ell_q}$.
\end{lemma}

\begin{proof}
\leanok
For finite sums this follows from \Cref{lem:lq-basic} by induction. The partial sums increase
to the series, and the $\ell_q$ norm of an increasing family of sequences is the supremum of
their norms, by monotone convergence. The truncations of $a$ increase to $a$, so the same
argument gives the last identity.
\end{proof}

The following two inequalities are the discrete Hardy inequalities: a tail sum and a partial
sum are controlled in a dyadically weighted $\ell_q$ norm.

\begin{lemma}[Hardy inequalities]
\label{lem:hardy}
\lean{ENNReal.hardy_tail_le, ENNReal.hardy_sum_le, ENNReal.lqNorm_shift_le',
  ENNReal.lqNorm_shiftRight_le, ENNReal.tsum_dyadicWeight_neg, ENNReal.tsum_le_lqNorm_dyadicWeight,
  ENNReal.tsum_dyadicWeight_neg_one, ENNReal.lqNorm_dyadicWeight_neg_one_le, ENNReal.tsum_ite_le_eq,
  ENNReal.tsum_ite_ge_eq}
\leanok
\uses{lem:countable-minkowski}
Let $\alpha \in \R$ and $1 \leq q \leq \infty$. For every sequence $b$ in $[0, \infty]$,
\[
  \Bigl\| \Bigl(2^{k\alpha} \sum_{i=0}^\infty b_{k+i}\Bigr)_k \Bigr\|_{\ell_q}
  \leq (1 - 2^{-\alpha})^{-1} \, \|(2^{k\alpha} b_k)_k\|_{\ell_q}
\]
and
\[
  \Bigl\| \Bigl(2^{-k\alpha} \sum_{j=0}^{k} b_j\Bigr)_k \Bigr\|_{\ell_q}
  \leq (1 - 2^{-\alpha})^{-1} \, \|(2^{-k\alpha} b_k)_k\|_{\ell_q} .
\]
In particular, $\sum_k b_k \leq (1 - 2^{-\alpha})^{-1} \, \|(2^{k\alpha} b_k)_k\|_{\ell_q}$.
\end{lemma}

\begin{proof}
\leanok
Write $2^{k\alpha} \sum_i b_{k+i} = \sum_i 2^{-i\alpha} \, 2^{(k+i)\alpha} b_{k+i}$. By
\Cref{lem:countable-minkowski} the left side is at most
$\sum_i 2^{-i\alpha} \|(2^{(k+i)\alpha} b_{k+i})_k\|_{\ell_q}$, a shift does not increase the
norm, and $\sum_i 2^{-i\alpha} = (1 - 2^{-\alpha})^{-1}$. The second inequality is the same
argument with $2^{-k\alpha} \sum_{j \leq k} b_j = \sum_{i \leq k} 2^{-i\alpha} \,
2^{-(k-i)\alpha} b_{k-i}$ and the shift to the right. The term $k = 0$ of the first inequality
gives the last one.
\end{proof}

The comparison of Pietsch's quasi-norm with its dyadic form rests on Cauchy's condensation
test with power weights.

\begin{lemma}[Cauchy condensation with power weights]
\label{lem:condensation}
\lean{ENNReal.tsum_pow_mul_le_dyadic, ENNReal.dyadic_le_tsum_pow_mul,
  ENNReal.iSup_pow_mul_le_dyadic, ENNReal.dyadic_le_iSup_pow_mul, ENNReal.sum_range_two_pow_eq,
  ENNReal.rpow_le_rpow_of_nonpos, ENNReal.natCast_two_pow_eq, ENNReal.card_Ico_two_pow,
  ENNReal.natCast_succ_rpow_le, ENNReal.dyadicWeight_le_natCast_succ_rpow,
  ENNReal.sum_Ico_pow_mul_le, ENNReal.le_sum_Ico_pow_mul}
\leanok
\uses{def:lq-norm}
Let $b$ be a decreasing sequence in $[0, \infty]$, let $\gamma \in \R$, and put
$\gamma_+ = \max(\gamma, 0)$ and $\gamma_- = \max(-\gamma, 0)$. Then
\[
  \sum_{n=0}^\infty (n+1)^\gamma b_n \leq b_0 + 2^{\gamma_+} \sum_{k=0}^\infty
    2^{k(\gamma+1)} b_{2^k}
\]
and
\[
  \sum_{k=0}^\infty 2^{k(\gamma+1)} b_{2^{k+1}} \leq 2^{\gamma_-} \sum_{n=0}^\infty
    (n+1)^\gamma b_n ,
\]
and for every $s \geq 0$,
\[
  \sup_n (n+1)^s b_n \leq b_0 + 2^s \sup_k 2^{ks} b_{2^k},
  \qquad
  \sup_k 2^{ks} b_{2^k} \leq \sup_n (n+1)^s b_n .
\]
\end{lemma}

\begin{proof}
\leanok
Split $\N \setminus \{0\}$ into the blocks $2^k \leq n < 2^{k+1}$ of $2^k$ elements each. On
the block $k$, $(n+1)^\gamma$ lies between $2^{-\gamma_-} 2^{k\gamma}$ and
$2^{\gamma_+} 2^{k\gamma}$, and $b_n$ lies between $b_{2^{k+1}}$ and $b_{2^k}$.
\end{proof}

Three sums over dyadic scales serve the vector-valued singular integrals of
\Cref{sec:dilation-families} and the Bessel kernels of \Cref{sec:bessel-kernels}.

\begin{lemma}[The dyadic sum of $\min(t, t^{-1})$]
\label{lem:dyadic-min-sum}
\lean{ENNReal.sum_range_two_pow_le, ENNReal.exists_one_lt_two_pow_mul,
  ENNReal.sum_range_min_two_pow_le, ENNReal.tsum_min_two_pow_add_le,
  ENNReal.tsum_min_two_pow_le, ENNReal.min_inv_sq_le}
\leanok
For every $a \in [0, \infty]$,
\[
  \sum_{j \in \N} \min\bigl(2^j a, (2^j a)^{-1}\bigr) \leq 4 ,
\]
and $\min(x, x^{-1})^2 \leq \min(x, x^{-1})$ for every $x \in [0, \infty]$.
\end{lemma}

\begin{proof}
\leanok
The second claim holds since $\min(x, x^{-1}) \leq 1$. For $a = 0$ every term is $0$.
Otherwise some $j$ has $2^j a > 1$ (Mathlib's
\texttt{ENNReal.\allowbreak exists\_\allowbreak nat\_\allowbreak mul\_\allowbreak gt} and
$m < 2^m$); let $j_0$ be the least such $j$. For
$j < j_0$ the terms are at most $2^j a$, and
$\sum_{j < j_0} 2^j a \leq 2^{j_0} a = 2 \cdot 2^{j_0 - 1} a \leq 2$ when $j_0 \geq 1$. For
$i \in \N$ the term of $j = i + j_0$ is at most
$(2^{i + j_0} a)^{-1} = 2^{-i} (2^{j_0} a)^{-1} \leq 2^{-i}$, and $\sum_i 2^{-i} = 2$.
\end{proof}

\begin{lemma}[Finitely many changed terms]
\label{lem:tsum-changed-terms}
\lean{ENNReal.tsum_le_tsum_add_sum_of_forall_ge_eq}
\leanok
For all sequences $(u_j)$, $(v_j)$ in $[0, \infty]$ and every $N \in \N$ with $u_j = v_j$ for
every $j \geq N$,
\[
  \sum_{j \in \N} u_j \leq \sum_{j \in \N} v_j + \sum_{i < N} u_i .
\]
\end{lemma}

\begin{proof}
\leanok
The sum of the $u_j$ is $\sum_{i < N} u_i + \sum_{i \geq N} v_i$.
\end{proof}

\begin{lemma}[Dyadic scales of an annulus]
\label{lem:dyadic-annulus-count}
\lean{Nat.le_add_of_dyadic_annulus, ENNReal.tsum_le_of_mem_Icc,
  ENNReal.tsum_le_of_dyadic_annulus}
\leanok
Let $a > 0$ and $t$ be real and $N \in \N$. For all $i, j \in \N$ with $a \leq t / 2^j$ and
$t / 2^i \leq 2^N a$, $j \leq i + N$. Let $M \in [0, \infty]$ and let $(u_j)_{j \in \N}$ be a
sequence in $[0, \infty]$ with $u_j \leq M$ for every $j$. If there is $i_0 \in \N$ such that
$u_j \neq 0$ implies $i_0 \leq j \leq i_0 + N$, or if $u_j \neq 0$ implies
$a \leq t / 2^j \leq 2^N a$, then $\sum_j u_j \leq (N + 1) M$.
\end{lemma}

\begin{proof}
\leanok
$2^j a \leq t \leq 2^{i+N} a$ gives the first claim. In the first case at most $N + 1$ terms are
nonzero. In the second, let $i_0$ be the least $j$ with $u_j \neq 0$; by the first claim every
$j$ with $u_j \neq 0$ satisfies $i_0 \leq j \leq i_0 + N$.
\end{proof}

\chapter{Approximation spaces}
\label{ch:approximation-spaces}

Throughout, $A$ is an approximation scheme in $X$. The following definition is
Pietsch's~\cite{Pietsch}, with his sequence $(n^{s-1/q} E_{n-1}(f))_{n \geq 1}$ indexed
from $0$.

\begin{definition}[Approximation space]
\label{def:approx-space}
\lean{ApproximationScheme.approxSeq, ApproximationScheme.approxNorm, ApproximationScheme.MemApprox,
  ApproximationScheme.approxNorm_top, ApproximationScheme.approxNorm_rpow,
  ApproximationScheme.memApprox_iff, ApproximationScheme.MemApprox.approxNorm_ne_top}
\leanok
\uses{def:error, def:lq-norm}
Let $s \in \R$ and $1 \leq q \leq \infty$. For $f \in X$,
\[
  \|f\|_{A^s_q} = \bigl\| \bigl((n+1)^{s - 1/q} E_n(f)\bigr)_{n \in \N} \bigr\|_{\ell_q}
  \in [0, \infty],
\]
with $1/\infty = 0$. The \emph{approximation space} $A^s_q$ consists of the $f \in X$ with
$\|f\|_{A^s_q} < \infty$.
\end{definition}

\begin{definition}[Dyadic quasi-norm]
\label{def:dyadic-norm}
\lean{ApproximationScheme.approxSeqDyadic, ApproximationScheme.approxNormDyadic,
  ApproximationScheme.approxSeqDyadic_le_lqNorm, ApproximationScheme.enorm_le_approxNormDyadic,
  ApproximationScheme.lqNorm_le_approxNormDyadic, ApproximationScheme.approxSeqDyadic_rpow,
  ApproximationScheme.approxSeqDyadic_zero, ApproximationScheme.approxNormDyadic_zero,
  ApproximationScheme.approxSeqDyadic_smul, ApproximationScheme.approxNormDyadic_smul,
  ApproximationScheme.one_le_two_rpow, ApproximationScheme.two_rpow_ne_top}
\leanok
\uses{def:error, def:lq-norm}
Let $s \in \R$ and $1 \leq q \leq \infty$. For $f \in X$,
\[
  \|f\|^{\mathrm{dy}}_{A^s_q} = \|f\| + \bigl\| \bigl(2^{ks} E_{2^k}(f)\bigr)_{k \in \N}
  \bigr\|_{\ell_q} .
\]
\end{definition}

The term $\|f\| = E_0(f)$ is needed because the dyadic sequence starts at $E_1$. The dyadic
form is the working norm: every later theorem is stated for it, and the next theorem
transports it to Pietsch's definition.

\begin{theorem}[Dyadic form of the quasi-norm]
\label{thm:dyadic-norm}
\lean{ApproximationScheme.approxNorm_le_approxNormDyadic,
  ApproximationScheme.approxNormDyadic_le_approxNorm,
  ApproximationScheme.memApprox_iff_approxNormDyadic_lt_top,
  ApproximationScheme.lqNorm_dyadic_succ_le,
  ApproximationScheme.memApprox_iff_lqNorm_dyadic_lt_top}
\leanok
\uses{def:approx-space, def:dyadic-norm, lem:condensation, lem:lq-basic}
Let $s \geq 0$ and $1 \leq q \leq \infty$. For every $f \in X$,
\[
  \|f\|_{A^s_q} \leq 2^s \, \|f\|^{\mathrm{dy}}_{A^s_q}
  \qquad \text{and} \qquad
  \|f\|^{\mathrm{dy}}_{A^s_q} \leq 4 \cdot 2^s \, \|f\|_{A^s_q} .
\]
In particular $f \in A^s_q$ if and only if $\|f\|^{\mathrm{dy}}_{A^s_q} < \infty$, that is, if and
only if $(2^{ks} E_{2^k}(f))_k \in \ell_q$.
\end{theorem}

\begin{proof}
\leanok
\uses{lem:error-basic}
For $q < \infty$ raise both sides to the power $q$: $\|f\|_{A^s_q}^q = \sum_n (n+1)^{sq-1}
E_n(f)^q$. The errors decrease, so \Cref{lem:condensation} with $\gamma = sq - 1$ and
$b_n = E_n(f)^q$ compares this series with $E_0(f)^q + \sum_k 2^{ksq} E_{2^k}(f)^q$ in both
directions, and $E_0(f) = \|f\|$. Taking $q$-th roots costs at most the factors stated,
using $(x + y)^{1/q} \leq x^{1/q} + y^{1/q}$ and, for the second bound,
$2^{(k+1)s} = 2^s 2^{ks}$ and $E_1(f) \leq \|f\|$. For $q = \infty$ use the supremum form of
\Cref{lem:condensation}.
\end{proof}

\begin{proposition}[Linear structure]
\label{prop:approx-linear}
\lean{ApproximationScheme.approxNorm_smul, ApproximationScheme.enorm_le_approxNorm,
  ApproximationScheme.approxNorm_eq_zero_iff, ApproximationScheme.approxNormDyadic_add_le,
  ApproximationScheme.approxNorm_add_le, ApproximationScheme.memApprox_of_mem,
  ApproximationScheme.approxSpace, ApproximationScheme.mem_approxSpace,
  ApproximationScheme.approxNorm_zero, ApproximationScheme.memApprox_zero,
  ApproximationScheme.MemApprox.smul, ApproximationScheme.MemApprox.neg,
  ApproximationScheme.approxNorm_anti_scheme, ApproximationScheme.MemApprox.of_subset_scheme,
  ApproximationScheme.approxSeqDyadic_add_le_succ, ApproximationScheme.MemApprox.add,
  ApproximationScheme.MemApprox.sub, ApproximationScheme.MemApprox.sum,
  ApproximationScheme.iUnion_subset_approxSpace}
\leanok
\uses{thm:dyadic-norm, lem:error-basic}
Let $s \geq 0$ and $1 \leq q \leq \infty$. For all $f, g \in X$, $c \in \R$ and $m \in \N$:
\begin{enumerate}
  \item $\|c f\|_{A^s_q} = |c| \, \|f\|_{A^s_q}$;
  \item $\|f\| \leq \|f\|_{A^s_q}$, and $\|f\|_{A^s_q} = 0$ only for $f = 0$;
  \item $\|f + g\|^{\mathrm{dy}}_{A^s_q} \leq 2 \cdot 2^s \bigl(\|f\|^{\mathrm{dy}}_{A^s_q} +
    \|g\|^{\mathrm{dy}}_{A^s_q}\bigr)$ and
    $\|f + g\|_{A^s_q} \leq 8 \cdot 8^s \bigl(\|f\|_{A^s_q} + \|g\|_{A^s_q}\bigr)$;
  \item $A_m \subseteq A^s_q$;
  \item a scheme with larger sets has smaller quasi-norms and larger approximation spaces.
\end{enumerate}
In particular $A^s_q$ is a linear subspace of $X$.
\end{proposition}

\begin{proof}
\leanok
The first item follows from $E_n(cf) = |c| E_n(f)$, the second from the term $n = 0$, which
is $E_0(f) = \|f\|$. For the third, $E_{2^{k+1}}(f + g) \leq E_{2^k}(f) + E_{2^k}(g)$ by
\Cref{lem:error-basic}, so $2^{(k+1)s} E_{2^{k+1}}(f+g) \leq 2^s (2^{ks} E_{2^k}(f) +
2^{ks} E_{2^k}(g))$; the term $k = 0$ is at most $\|f + g\|$, and Minkowski's inequality
gives the dyadic bound. Pietsch's form follows by \Cref{thm:dyadic-norm}. For the fourth,
$E_n(a) = 0$ for $a \in A_m$ and $n \geq m$, so only finitely many terms are nonzero.
\end{proof}

The constants are not optimised. The quasi-triangle inequality is all that is needed: Mathlib
has no class of quasi-normed spaces, so $A^s_q$ is a linear subspace of $X$ together with a
function into $[0, \infty]$, in the way Mathlib treats $L_p$ through \texttt{eLpNorm}.

\begin{lemma}[Approximation spaces lie in the closure]
\label{lem:approx-closure}
\lean{ApproximationScheme.approxSeq_pow_le,
  ApproximationScheme.tendsto_approxError_pow_zero_of_memApprox,
  ApproximationScheme.mem_closure_iUnion_of_memApprox,
  ApproximationScheme.approxSpace_subset_closure}
\leanok
\uses{def:dyadic-norm, lem:error-tendsto}
Let $s > 0$ and $1 \leq q \leq \infty$. For every $f \in A^s_q$ and $k \in \N$,
$E_{2^k}(f) \leq 2^{-ks} \|f\|^{\mathrm{dy}}_{A^s_q}$, and $f$ lies in the closure of
$\bigcup_n A_n$.
\end{lemma}

\begin{proof}
\leanok
\uses{thm:dyadic-norm}
The bound is $2^{ks} E_{2^k}(f) \leq \|(2^{js} E_{2^j}(f))_j\|_{\ell_q}$. It tends to zero,
and \Cref{lem:error-tendsto} applies.
\end{proof}

\begin{theorem}[Embeddings between approximation spaces]
\label{thm:approx-embedding}
\lean{ApproximationScheme.approxNormDyadic_mono_exponent,
  ApproximationScheme.approxNormDyadic_le_of_lt, ApproximationScheme.MemApprox.mono,
  ApproximationScheme.approxSpace_mono, ApproximationScheme.MemApprox.lqNorm_one_lt_top,
  ApproximationScheme.approxSeqDyadic_eq_mul, ApproximationScheme.approxNorm_le_of_exponent_le,
  ApproximationScheme.approxNorm_le_of_lt, ApproximationScheme.MemApprox.mono_exponent,
  ApproximationScheme.MemApprox.of_lt}
\leanok
\uses{def:dyadic-norm, lem:lq-basic, thm:dyadic-norm}
Let $1 \leq q, r \leq \infty$, let $s, t \in \R$ and $f \in X$.
\begin{enumerate}
  \item If $q \leq r$, then $\|f\|^{\mathrm{dy}}_{A^s_r} \leq \|f\|^{\mathrm{dy}}_{A^s_q}$.
  \item If $t < s$, then $\|f\|^{\mathrm{dy}}_{A^t_r} \leq (1 - 2^{-(s-t)})^{-1}
    \|f\|^{\mathrm{dy}}_{A^s_q}$.
  \item If $0 \leq t$, and $t < s$ or $t = s$ and $q \leq r$, then $A^s_q \subseteq A^t_r$.
  \item For Pietsch's quasi-norm: if $s \geq 0$ and $q \leq r$, then
    $\|f\|_{A^s_r} \leq 2^s \cdot 4 \cdot 2^s \, \|f\|_{A^s_q}$; if $0 \leq t < s$, then
    $\|f\|_{A^t_r} \leq 2^t (1 - 2^{-(s-t)})^{-1} \cdot 4 \cdot 2^s \, \|f\|_{A^s_q}$.
\end{enumerate}
\end{theorem}

\begin{proof}
\leanok
The first item is the monotonicity of the $\ell_q$ norms in $q$, since the dyadic weights do
not depend on $q$; Pietsch's weights do, which is why the dyadic form is used. For the
second, $2^{kt} E_{2^k}(f) = 2^{-k(s-t)} \, 2^{ks} E_{2^k}(f) \leq 2^{-k(s-t)}
\|(2^{js} E_{2^j}(f))_j\|_{\ell_q}$, and $\|(2^{-k(s-t)})_k\|_{\ell_r} \leq (1 -
2^{-(s-t)})^{-1}$. The third follows by \Cref{thm:dyadic-norm}, and the fourth from the first
two and the two comparisons of \Cref{thm:dyadic-norm}.
\end{proof}

\begin{lemma}[Functoriality]
\label{lem:approx-functorial}
\lean{ApproximationScheme.approxSeq_map_le, ApproximationScheme.approxNorm_map_le,
  ApproximationScheme.approxNormDyadic_map_le, ApproximationScheme.MemApprox.map,
  ApproximationScheme.approxSpace_map_le, ApproximationScheme.mem_map_of_mem'}
\leanok
\uses{lem:error-map, def:approx-space, def:dyadic-norm}
Let $A$ and $B$ be approximation schemes in $X$ and $Y$, and let $T : X \to Y$ be a bounded
linear map with $T(A_n) \subseteq B_n$ for all $n$. Then for all $s \in \R$,
$1 \leq q \leq \infty$ and $f \in X$,
\[
  \|Tf\|_{B^s_q} \leq \|T\| \, \|f\|_{A^s_q}
  \qquad \text{and} \qquad
  \|Tf\|^{\mathrm{dy}}_{B^s_q} \leq \|T\| \, \|f\|^{\mathrm{dy}}_{A^s_q} .
\]
\end{lemma}

\begin{proof}
\leanok
Apply \Cref{lem:error-map} to every term, and $\|Tf\| \leq \|T\| \|f\|$.
\end{proof}

\chapter{Lifting and the Nikolskii transfer theorem}
\label{ch:lifting}

Throughout, $A$ is an approximation scheme in $X$, and $\iota : Y \to X$ is a bounded linear
map. One hypothesis covers two classical inequalities. If $\iota$ is the embedding of a
space of smooth functions, it is a Bernstein inequality; if $\iota$ is the inclusion
$L_r \hookrightarrow L_p$, it is a Nikolskii inequality.

\begin{definition}[Lifting bound]
\label{def:lift-bound}
\lean{ApproximationScheme.LiftBound, ApproximationScheme.LiftBound.exists_enorm_le}
\leanok
\uses{def:scheme}
Let $\beta, C \in \R$. The map $\iota$ satisfies the \emph{lifting bound} of exponent $\beta$
and constant $C$ for $A$ if for every $n \geq 1$ and every $a \in A_n$ there is $y \in Y$ with
$\iota y = a$ and $\|y\| \leq C \, n^\beta \, \|a\|$.
\end{definition}

\begin{lemma}[Near-best dyadic approximants]
\label{lem:near-best}
\lean{ApproximationScheme.DyadicApprox, ApproximationScheme.exists_dyadicApprox,
  ApproximationScheme.DyadicApprox.sub_mem, ApproximationScheme.DyadicApprox.enorm_sub_succ_le,
  ApproximationScheme.DyadicApprox.tendsto, ApproximationScheme.DyadicApprox.enorm_zero_le}
\leanok
\uses{lem:error-basic, lem:scheme-basic}
Let $f \in X$ and let $(t_k)_k$ be a sequence in $(0, \infty)$. There are $a_k \in A_{2^k}$
with $\|f - a_k\| \leq E_{2^k}(f) + t_k$ for every $k$. Every such sequence satisfies
$a_{k+1} - a_k \in A_{2^{k+2}}$; if $(t_k)$ decreases, then
$\|a_{k+1} - a_k\| \leq 2 (E_{2^k}(f) + t_k)$; and if $E_{2^k}(f) \to 0$ and $t_k \to 0$, then
$a_k \to f$.
\end{lemma}

\begin{proof}
\leanok
The approximants exist by the definition of the infimum. The rest is
\Cref{lem:scheme-basic}(5), the triangle inequality and the monotonicity of the errors.
\end{proof}

The tolerances $t_k$ replace best approximants, which need not exist. The proofs below take
$t_k = \|f\| \, 2^{-k(\sigma+1)}$ for a suitable $\sigma$, so that
$\sum_k 2^{k\sigma} t_k \leq 2 \|f\|$.

\begin{theorem}[Lifting theorem]
\label{thm:lift}
\lean{ApproximationScheme.exists_preimage_of_liftBound, ApproximationScheme.liftConst,
  ApproximationScheme.DyadicApprox.exists_lift, ApproximationScheme.DyadicApprox.exists_preimage,
  ApproximationScheme.tolerance, ApproximationScheme.liftConst_ne_top,
  ApproximationScheme.DyadicApprox.exists_lift_partial, ApproximationScheme.tolerance_ne_zero,
  ApproximationScheme.tolerance_ne_top, ApproximationScheme.tolerance_anti,
  ApproximationScheme.tolerance_zero, ApproximationScheme.dyadicWeight_mul_tolerance_le,
  ApproximationScheme.tsum_dyadicWeight_mul_tolerance_le, ApproximationScheme.tendsto_tolerance,
  ApproximationScheme.tsum_tolerance_le, ApproximationScheme.lqNorm_dyadicWeight_mul_tolerance_le}
\leanok
\uses{def:lift-bound, lem:near-best, def:dyadic-norm}
Let $Y$ be a Banach space, let $\beta \geq 0$ and $C \geq 0$, and let $\iota$ satisfy the
lifting bound of exponent $\beta$ and constant $C$ for $A$. For every $f \in X$ with
$\sum_k 2^{k\beta} E_{2^k}(f) < \infty$ there is $y \in Y$ with $\iota y = f$ and
\[
  \|y\| \leq (3C + 6C \cdot 4^\beta) \Bigl(\|f\| + \sum_{k=0}^\infty 2^{k\beta}
    E_{2^k}(f)\Bigr) .
\]
\end{theorem}

\begin{proof}
\leanok
For $f = 0$ take $y = 0$. Otherwise choose near-best approximants $a_k$ with tolerances
$t_k = \|f\| 2^{-k(\beta+1)}$ by \Cref{lem:near-best}. Lift $a_0 \in A_1$ to $y_0$ and
$a_{k+1} - a_k \in A_{2^{k+2}}$ to $d_k$, so that $\|y_0\| \leq C \|a_0\| \leq 3C\|f\|$
and $\|d_k\| \leq C \, 2^{(k+2)\beta} \cdot 2(E_{2^k}(f) + t_k) = 2C \, 4^\beta \,
2^{k\beta} (E_{2^k}(f) + t_k)$. The series $\sum_k \|d_k\|$ converges, so
$y = y_0 + \sum_k d_k$ exists in $Y$. By continuity of $\iota$, $\iota y$ is the limit of
$\iota(y_0 + \sum_{j<k} d_j) = a_k$, which is $f$ since $E_{2^k}(f) \to 0$. Adding the bounds
gives the constant.
\end{proof}

The same construction gives the next theorem, the transfer of approximation spaces along an
injective map with a Nikolskii inequality.

\begin{theorem}[Nikolskii transfer theorem]
\label{thm:nikolskii-transfer}
\lean{ApproximationScheme.exists_preimage_of_nikolskii,
  ApproximationScheme.MemApprox.exists_preimage_of_nikolskii, ApproximationScheme.nikolskiiConst,
  ApproximationScheme.nikolskiiConst_ne_top}
\leanok
\uses{thm:lift, lem:hardy, thm:approx-embedding, def:comap}
Let $Y$ be a Banach space, let $\iota$ be injective, let $C \geq 0$ and $0 \leq \alpha < s$,
and let $\iota$ satisfy the lifting bound of exponent $\alpha$ and constant $C$ for $A$. Let
$B$ be the scheme $B_n = \iota^{-1}(A_n)$ in $Y$, and let $1 \leq q \leq \infty$. For every
$f \in A^s_q$ there is $y \in Y$ with $\iota y = f$ and
\[
  \|y\|^{\mathrm{dy}}_{B^{s-\alpha}_q} \leq K \, \|f\|^{\mathrm{dy}}_{A^s_q},
  \qquad
  K = 3C + 6C \cdot 4^\alpha (1 - 2^{-(s-\alpha)})^{-1} + 4C \cdot 4^\alpha .
\]
\end{theorem}

\begin{proof}
\leanok
By \Cref{thm:approx-embedding}, $\sum_k 2^{k\alpha} E_{2^k}(f) < \infty$, so the
construction of \Cref{thm:lift} applies, now with tolerances $t_k = \|f\| 2^{-k(s+1)}$. The
partial lifts $z_k = y_0 + \sum_{j<k} d_j$ satisfy $\iota z_k = a_k \in A_{2^k}$, so
$z_k \in B_{2^k}$ and
\[
  E^B_{2^k}(y) \leq \|y - z_k\| \leq \sum_{i=0}^\infty \|d_{k+i}\|
  \leq 2C \, 4^\alpha \sum_{i=0}^\infty 2^{(k+i)\alpha} \bigl(E_{2^{k+i}}(f) + t_{k+i}\bigr).
\]
Multiply by $2^{k(s-\alpha)}$ and apply the first Hardy inequality (\Cref{lem:hardy}) with
exponent $s - \alpha > 0$: the weights combine to $2^{k(s-\alpha)} 2^{k\alpha} = 2^{ks}$,
and $\|(2^{ks} t_k)_k\|_{\ell_q} \leq 2\|f\|$.
\end{proof}

\section{Jackson and Bernstein inequalities}

\begin{definition}[Jackson inequality]
\label{def:jackson}
\lean{ApproximationScheme.JacksonBound, ApproximationScheme.JacksonBound.approxSeq_le}
\leanok
\uses{def:error}
Let $r, C \in \R$. The map $\iota$ satisfies \emph{Jackson's inequality} of order $r$ and
constant $C$ for $A$ if $E_n(\iota y) \leq C \, n^{-r} \|y\|$ for every $n \geq 1$ and every
$y \in Y$.
\end{definition}

\begin{lemma}[Jackson's inequality gives $Y \subseteq A^r_\infty$]
\label{lem:jackson-top}
\lean{ApproximationScheme.approxNormDyadic_top_le_of_jackson,
  ApproximationScheme.memApprox_top_of_jackson}
\leanok
\uses{def:jackson, def:dyadic-norm, thm:dyadic-norm}
Let $\iota$ satisfy Jackson's inequality of order $r \geq 0$ and constant $C \geq 0$. Then
$\|\iota y\|^{\mathrm{dy}}_{A^r_\infty} \leq (\|\iota\| + C) \|y\|$ for every $y \in Y$; in
particular $\iota y \in A^r_\infty$.
\end{lemma}

\begin{proof}
\leanok
$2^{kr} E_{2^k}(\iota y) \leq C \|y\|$ for every $k$, and $\|\iota y\| \leq \|\iota\| \|y\|$.
\end{proof}

\begin{definition}[Smoothness space]
\label{def:smoothness}
\lean{ApproximationScheme.SmoothnessSpace}
\leanok
\uses{def:jackson, def:lift-bound}
A \emph{smoothness space} of order $r > 0$ for $A$ is a normed space $Y$ with an injective
bounded linear map $\iota : Y \to X$ that satisfies Jackson's inequality of order $r$ and the
lifting bound of exponent $r$ for $A$, both with nonnegative constants.
\end{definition}

The lifting bound is then Bernstein's inequality $\|\iota^{-1} a\|_Y \leq C n^r \|a\|$ for
$a \in A_n$. Sobolev spaces of order $r$ inside $L_p$ are the motivating example.

\begin{proposition}[Smoothness spaces between $A^r_1$ and $A^r_\infty$]
\label{prop:smoothness}
\lean{ApproximationScheme.SmoothnessSpace.memApprox_top,
  ApproximationScheme.SmoothnessSpace.exists_eq_of_memApprox_one,
  ApproximationScheme.SmoothnessSpace.exists_eq_of_lt, ApproximationScheme.SmoothnessSpace.congr}
\leanok
\uses{def:smoothness, lem:jackson-top, thm:lift, thm:approx-embedding}
Let $(Y, \iota)$ be a smoothness space of order $r$ for $A$. Then $\iota(Y) \subseteq
A^r_\infty$. If $Y$ is a Banach space, then $A^r_1 \subseteq \iota(Y)$, and
$A^s_q \subseteq \iota(Y)$ for every $s > r$ and $1 \leq q \leq \infty$.
\end{proposition}

\begin{proof}
\leanok
The first inclusion is \Cref{lem:jackson-top}, the second \Cref{thm:lift} with $\beta = r$.
The third follows from the second and $A^s_q \subseteq A^r_1$ (\Cref{thm:approx-embedding}).
\end{proof}

\chapter{Interpolation and the representation theorem}
\label{ch:interpolation}

Throughout, $\iota : Y \to X$ is a bounded linear map. The $K$-functional of Peetre measures
how well $f \in X$ is approximated by elements of $Y$ of controlled norm; the real
interpolation spaces are defined from it~\cite{BerghLofstrom}.

\begin{definition}[$K$-functional]
\label{def:k-functional}
\lean{ApproximationScheme.kFunctional, ApproximationScheme.kFunctional_le}
\leanok
For $t \in [0, \infty]$ and $f \in X$,
\[
  K(t, f) = \inf_{y \in Y} \bigl(\|f - \iota y\| + t \|y\|\bigr) .
\]
\end{definition}

\begin{lemma}[Properties of the $K$-functional]
\label{lem:k-functional}
\lean{ApproximationScheme.kFunctional_le_enorm, ApproximationScheme.kFunctional_mono,
  ApproximationScheme.kFunctional_add_le, ApproximationScheme.kFunctional_smul_le,
  ApproximationScheme.kFunctional_zero_right}
\leanok
\uses{def:k-functional}
For all $t \leq u$ in $[0, \infty]$, $f, g \in X$ and $c \in \R$: $K(t, f) \leq \|f\|$,
$K(t, f) \leq K(u, f)$, $K(t, f + g) \leq K(t, f) + K(t, g)$ and
$K(t, c f) \leq |c| \, K(t, f)$.
\end{lemma}

\begin{proof}
\leanok
Take $y = 0$ for the first bound; the others hold for each $y$ before the infimum.
\end{proof}

\begin{definition}[Real interpolation space]
\label{def:interpolation}
\lean{ApproximationScheme.interpSeq, ApproximationScheme.interpNorm, ApproximationScheme.MemInterp,
  ApproximationScheme.interpSpace, ApproximationScheme.interpNorm_add_le,
  ApproximationScheme.enorm_le_interpNorm, ApproximationScheme.interpNorm_smul_le,
  ApproximationScheme.mem_interpSpace}
\leanok
\uses{def:k-functional, def:lq-norm, lem:k-functional}
Let $\theta, \rho \in \R$ and $1 \leq q \leq \infty$. For $f \in X$,
\[
  \|f\|_{\theta, q} = \|f\| + \bigl\| \bigl(2^{k\rho\theta} K(2^{-k\rho}, f)\bigr)_{k \in \N}
  \bigr\|_{\ell_q} .
\]
The \emph{real interpolation space} $(X, Y)_{\theta, q}$ consists of the $f \in X$ with
$\|f\|_{\theta, q} < \infty$. It is a linear subspace of $X$, and $\|\cdot\|_{\theta, q}$
satisfies the triangle inequality.
\end{definition}

This is the discrete $K$-method. The textbook definition samples $K(t, f)$ on the grid
$t = 2^{-k}$, which is $\rho = 1$; the representation theorem produces the grid $t = 2^{-kr}$.
Grids for different $\rho > 0$ give equivalent norms, since $K(t, f)$ increases and $K(t, f)/t$
decreases in $t$, and the continuous $K$-method
$\bigl(\int_0^\infty (t^{-\theta} K(t, f))^q \, dt/t\bigr)^{1/q}$ gives an equivalent norm as
well; these equivalences are not formalised. The term $\|f\|$ replaces the negative indices
$k < 0$ of the textbook definition.

\begin{lemma}[Errors and the $K$-functional]
\label{lem:error-k-functional}
\lean{ApproximationScheme.approxSeq_le_kFunctional_of_jackson}
\leanok
\uses{def:k-functional, def:jackson}
Let $A$ be an approximation scheme in $X$, and let $\iota$ satisfy Jackson's inequality of
order $r$ and constant $C \geq 0$. For every $n \geq 1$ and $f \in X$,
$E_n(f) \leq \max(1, C) \, K(n^{-r}, f)$.
\end{lemma}

\begin{proof}
\leanok
\uses{lem:error-basic}
For every $y \in Y$, $E_n(f) \leq \|f - \iota y\| + E_n(\iota y) \leq \|f - \iota y\| + C
n^{-r} \|y\|$.
\end{proof}

The representation theorem identifies approximation spaces with interpolation spaces, as in
DeVore and Lorentz~\cite[Chapter~7]{DeVoreLorentz}.

\begin{theorem}[Representation theorem]
\label{thm:representation}
\lean{ApproximationScheme.approxNormDyadic_le_interpNorm_of_jackson,
  ApproximationScheme.interpSeq_le_of_lift,
  ApproximationScheme.interpNorm_le_approxNormDyadic_of_bernstein, ApproximationScheme.reprConst,
  ApproximationScheme.memApprox_iff_memInterp, ApproximationScheme.approxSpace_eq_interpSpace,
  ApproximationScheme.SmoothnessSpace.memApprox_iff_memInterp, ApproximationScheme.reprConst_ne_top}
\leanok
\uses{lem:error-k-functional, def:interpolation, lem:near-best, lem:hardy,
  def:lift-bound, def:dyadic-norm}
Let $A$ be an approximation scheme in $X$, let $r \in \R$ and $C_1, C_2 \geq 0$, and let
$\iota$ satisfy Jackson's inequality of order $r$ and constant $C_1$ and the lifting bound of
exponent $r$ and constant $C_2$. Let $0 < s < r$ and $1 \leq q \leq \infty$. Then for every
$f \in X$,
\[
  \|f\|^{\mathrm{dy}}_{A^s_q} \leq \max(1, C_1) \, \|f\|_{s/r, q}
  \qquad \text{and} \qquad
  \|f\|_{s/r, q} \leq K \, \|f\|^{\mathrm{dy}}_{A^s_q},
\]
where $K = 3 + (3C_2 + 4C_2 \cdot 4^r)(1 - 2^{-(r-s)})^{-1}$ and the interpolation norm is
taken on the grid $\rho = r$. In particular $A^s_q = (X, Y)_{s/r, q}$.
\end{theorem}

\begin{proof}
\leanok
For the first inequality apply \Cref{lem:error-k-functional} with $n = 2^k$: then
$n^{-r} = 2^{-kr}$ and $2^{ks} = 2^{kr \cdot s/r}$. For the second choose near-best
approximants $a_k$ (\Cref{lem:near-best}) and lift them as in \Cref{thm:lift}: the partial
lifts $z_k = y_0 + \sum_{j<k} d_j$ satisfy $\iota z_k = a_k$ and
$\|z_k\| \leq 3C_2 \|f\| + 2C_2 4^r \sum_{j<k} 2^{jr} (E_{2^j}(f) + t_j)$. Hence
\[
  2^{ks} K(2^{-kr}, f) \leq 2^{ks} \|f - a_k\| + 2^{-k(r-s)} \|z_k\| ,
\]
and the second Hardy inequality (\Cref{lem:hardy}) with exponent $r - s > 0$ bounds the
partial sums. The tolerances are made arbitrarily small at the end. Only finitely many
lifts enter each $z_k$, so $Y$ need not be complete.
\end{proof}

Here $(X, Y)_{s/r, q}$ is the discrete $K$-method on the grid $t = 2^{-kr}$
(\Cref{def:interpolation}); its equivalence with the continuous $K$-method is not formalised.

\begin{corollary}[Schemes with the same Jackson and Bernstein inequalities]
\label{cor:equal-spaces}
\lean{ApproximationScheme.approxSpace_eq_of_jackson_bernstein}
\leanok
\uses{thm:representation}
Let $A$ and $B$ be approximation schemes in $X$, let $r \in \R$, and let $\iota$ satisfy
Jackson's inequality of order $r$ and the lifting bound of exponent $r$ for $A$ and for $B$,
with nonnegative constants. Then $A^s_q = B^s_q$ for all $0 < s < r$ and
$1 \leq q \leq \infty$.
\end{corollary}

\begin{proof}
\leanok
Both are $(X, Y)_{s/r, q}$ by \Cref{thm:representation}.
\end{proof}

This is how approximation by trigonometric polynomials, splines and wavelets are seen to
define the same spaces.

\chapter{Block decompositions and projections}
\label{ch:decompositions}

Throughout, $A$ is an approximation scheme in $X$. Approximation spaces can be described by
block decompositions $f = \sum_k g_k$ into blocks $g_k \in A_{2^k}$~\cite[Chapter~7]{DeVoreLorentz}.
Each block is a whole level; atomic decompositions refine the blocks into localised atoms with
coefficients, and wavelet decompositions expand them in a basis.

\begin{definition}[Block decomposition norm]
\label{def:block-decomposition}
\lean{ApproximationScheme.blockDecompNorm, ApproximationScheme.sum_range_two_pow_lt,
  ApproximationScheme.blockDecompNorm_le, ApproximationScheme.sum_range_mem}
\leanok
\uses{def:scheme, def:lq-norm}
Let $s \in \R$ and $1 \leq q \leq \infty$. For $f \in X$, $\|f\|^{\mathrm{dec}}_{s,q}$ is
the infimum of $\|(2^{ks} \|g_k\|)_k\|_{\ell_q}$ over all sequences $(g_k)_k$ with
$g_k \in A_{2^k}$ for every $k$ and $\sum_k g_k = f$, the series converging in $X$.
\end{definition}

\begin{theorem}[Block decompositions]
\label{thm:block-decomposition}
\lean{ApproximationScheme.approxNormDyadic_le_of_hasSum,
  ApproximationScheme.approxNormDyadic_le_blockDecompNorm,
  ApproximationScheme.blockDecompNorm_le_approxNormDyadic,
  ApproximationScheme.memApprox_iff_blockDecompNorm_lt_top,
  ApproximationScheme.DyadicApprox.blocks}
\leanok
\uses{def:block-decomposition, def:dyadic-norm, lem:hardy, lem:near-best, lem:approx-closure}
Let $s > 0$ and $1 \leq q \leq \infty$. For every $f \in X$,
\[
  \|f\|^{\mathrm{dy}}_{A^s_q} \leq 2 (1 - 2^{-s})^{-1} \, \|f\|^{\mathrm{dec}}_{s,q} .
\]
If $X$ is a Banach space, then for every $f \in X$,
\[
  \|f\|^{\mathrm{dec}}_{s,q} \leq (3 \cdot 2^s + 4 \cdot 4^s) \, \|f\|^{\mathrm{dy}}_{A^s_q},
\]
and $f \in A^s_q$ if and only if $\|f\|^{\mathrm{dec}}_{s,q} < \infty$.
\end{theorem}

\begin{proof}
\leanok
\uses{lem:scheme-basic}
Let $f = \sum_k g_k$ with $g_k \in A_{2^k}$. Since $\sum_{j<k} 2^j < 2^k$, the partial sum
$\sum_{j<k} g_j$ lies in $A_{2^k}$, so $E_{2^k}(f) \leq T_k = \sum_{i \geq 0} \|g_{k+i}\|$,
and also $\|f\| \leq T_0$. The first Hardy inequality bounds $\|(2^{ks} T_k)_k\|_{\ell_q}$,
which gives the first inequality. Conversely, take near-best approximants $a_k$ with
tolerances $\|f\| 2^{-k(s+1)}$ and the blocks $g_0 = 0$, $g_1 = a_0$, $g_{k+2} =
a_{k+1} - a_k \in A_{2^{k+2}}$. Their partial sums are the $a_k$, which converge to $f$ by
\Cref{lem:approx-closure}, and $\|g_{k+2}\| \leq 2(E_{2^k}(f) + t_k)$ gives the bound.
\end{proof}

\section{Projection schemes}

\begin{definition}[Projection scheme]
\label{def:projection-scheme}
\lean{ApproximationScheme.ProjectionScheme, ApproximationScheme.ProjectionScheme.bound_nonneg,
  ApproximationScheme.ProjectionScheme.enorm_apply_le}
\leanok
\uses{def:scheme}
A \emph{projection scheme} for $A$ is a sequence of bounded linear maps $P_n : X \to X$ with
$P_n(X) \subseteq A_n$, $P_n a = a$ for every $a \in A_n$, and $\|P_n\| \leq M$ for all $n$,
for some $M \geq 0$.
\end{definition}

\begin{lemma}[Lebesgue inequality]
\label{lem:lebesgue}
\lean{ApproximationScheme.ProjectionScheme.approxSeq_le_enorm_sub,
  ApproximationScheme.ProjectionScheme.enorm_sub_le_approxSeq,
  ApproximationScheme.ProjectionScheme.tendsto_of_dense}
\leanok
\uses{def:projection-scheme, def:error, lem:error-tendsto}
Let $(P_n)$ be a projection scheme for $A$ with bound $M$. For every $n \in \N$ and
$f \in X$,
\[
  E_n(f) \leq \|f - P_n f\| \leq (1 + M) \, E_n(f) .
\]
If $\bigcup_n A_n$ is dense in $X$, then $P_n f \to f$ for every $f \in X$.
\end{lemma}

\begin{proof}
\leanok
For $a \in A_n$, $f - P_n f = (f - a) - P_n(f - a)$. The last statement follows from
$E_n(f) \to 0$.
\end{proof}

\begin{definition}[Dyadic blocks]
\label{def:blocks}
\lean{ApproximationScheme.ProjectionScheme.block,
  ApproximationScheme.ProjectionScheme.blockNorm}
\leanok
\uses{def:projection-scheme, def:lq-norm}
For a projection scheme $(P_n)$, $k \in \N$ and $f \in X$, the \emph{dyadic block} is
$\Delta_k f = P_{2^{k+1}} f - P_{2^k} f$. For $s \in \R$ and $1 \leq q \leq \infty$,
\[
  \|f\|^{\mathrm{bl}}_{s,q} = \|P_1 f\| + \bigl\| \bigl(2^{ks} \|\Delta_k f\|\bigr)_k
  \bigr\|_{\ell_q} .
\]
\end{definition}

\begin{theorem}[Characterisation by dyadic blocks]
\label{thm:blocks}
\lean{ApproximationScheme.ProjectionScheme.enorm_block_le,
  ApproximationScheme.ProjectionScheme.blockNorm_le_approxNormDyadic,
  ApproximationScheme.ProjectionScheme.enorm_sub_le_tsum_block,
  ApproximationScheme.ProjectionScheme.approxNormDyadic_le_blockNorm,
  ApproximationScheme.ProjectionScheme.memApprox_iff_blockNorm_lt_top,
  ApproximationScheme.ProjectionScheme.sum_block}
\leanok
\uses{def:blocks, lem:lebesgue, lem:hardy, def:dyadic-norm}
Let $(P_n)$ be a projection scheme for $A$ with bound $M$, let $s \in \R$ and
$1 \leq q \leq \infty$. For every $f \in X$,
\[
  \|f\|^{\mathrm{bl}}_{s,q} \leq (M + 2(1 + M)) \, \|f\|^{\mathrm{dy}}_{A^s_q} .
\]
If $P_n f \to f$, then
\[
  \|f\|^{\mathrm{dy}}_{A^s_q} \leq \bigl(1 + 2(1 - 2^{-s})^{-1}\bigr) \,
  \|f\|^{\mathrm{bl}}_{s,q} .
\]
In particular, if $\bigcup_n A_n$ is dense in $X$ and $s > 0$, then $f \in A^s_q$ if and
only if $\|f\|^{\mathrm{bl}}_{s,q} < \infty$.
\end{theorem}

\begin{proof}
\leanok
By \Cref{lem:lebesgue}, $\|\Delta_k f\| \leq \|f - P_{2^{k+1}} f\| + \|f - P_{2^k} f\| \leq
2(1 + M) E_{2^k}(f)$, and $\|P_1 f\| \leq M \|f\|$. Conversely, if $P_n f \to f$, then
$f - P_{2^k} f = \sum_{i \geq 0} \Delta_{k+i} f$, so $E_{2^k}(f) \leq \sum_{i \geq 0}
\|\Delta_{k+i} f\|$ and $\|f\| \leq \|P_1 f\| + \sum_{i \geq 0} \|\Delta_i f\|$; the first
Hardy inequality concludes.
\end{proof}

\section{Synthesis with decay and atoms}

Atoms of a function space lie in no fixed set of approximants. Their blocks satisfy a weaker
condition: they are small and approximable at a rate $\sigma$ from their own level on.

\begin{definition}[Linear scheme]
\label{def:linear-scheme}
\lean{ApproximationScheme.IsLinear, ApproximationScheme.isLinear_ofSubmodules,
  ApproximationScheme.isLinear_ofMonotoneSubmodules,
  ApproximationScheme.approxSeq_add_le_of_isLinear,
  ApproximationScheme.approxSeq_sum_le_of_isLinear, ApproximationScheme.IsLinear.map,
  ApproximationScheme.approxError_add_le_of_isLinear}
\leanok
\uses{def:scheme, def:error}
An approximation scheme is \emph{linear} if $a + b \in A_n$ for all $n \in \N$ and
$a, b \in A_n$. For a linear scheme, $E_n(f + g) \leq E_n(f) + E_n(g)$ for all $n$, $f$, $g$.
\end{definition}

\begin{theorem}[Synthesis with decay]
\label{thm:synthesis-decay}
\lean{ApproximationScheme.approxNormDyadic_le_of_hasSum_decay,
  ApproximationScheme.dyadicWeight_mul_dyadicWeight_sub}
\leanok
\uses{def:linear-scheme, def:dyadic-norm, lem:hardy, lem:error-basic}
Let $A$ be a linear scheme, $\sigma \in \R$, $C \in [0, \infty]$ and $1 \leq q \leq \infty$.
Let $f = \sum_j g_j$ in $X$, and let $(\varepsilon_j)_j$ in $[0, \infty]$ satisfy
$\|g_j\| \leq \varepsilon_j$ for every $j$ and $E_{2^k}(g_j) \leq C \, 2^{-(k-j)\sigma}
\varepsilon_j$ for all $j \leq k$. Then for every $s \in \R$,
\[
  \|f\|^{\mathrm{dy}}_{A^s_q} \leq \bigl(2 (1 - 2^{-s})^{-1} + C (1 - 2^{-(\sigma - s)})^{-1}\bigr)
  \bigl\|(2^{js} \varepsilon_j)_j\bigr\|_{\ell_q} ,
\]
a finite constant for $0 < s < \sigma$.
\end{theorem}

\begin{proof}
\leanok
Split $f$ at level $k$: $E_{2^k}(f) \leq \sum_{j \leq k} E_{2^k}(g_j) + \|\sum_{j > k} g_j\|$
by linearity. The first sum is at most $C \, 2^{-k\sigma} \sum_{j \leq k} 2^{j\sigma}
\varepsilon_j$, and after multiplication by $2^{ks}$ the second Hardy inequality with exponent
$\sigma - s$ applies. The second term is at most $\sum_{i \geq 0} \varepsilon_{k+i}$, and the
first Hardy inequality with exponent $s$ applies; it also bounds $\|f\| \leq \sum_j
\varepsilon_j$.
\end{proof}

This is the pattern of the proof of \Cref{thm:moduli-coincidence}. For atoms the size
$\varepsilon_j$ is that of a whole level, measured by the coefficients.

\begin{definition}[Atomic system and the sequence space $b^s_{p,q}$]
\label{def:atomic-system}
\lean{ApproximationScheme.AtomicSystem, ApproximationScheme.bNorm}
\leanok
\uses{def:scheme, def:error, def:lq-norm}
Let $M_j$ be finite sets, $1 \leq p \leq \infty$ and $\delta, \sigma \in \R$. An \emph{atomic
system} for $A$ consists of sets $\mathcal{A}_{j,m} \subseteq X$ of admissible atoms, $j \in
\N$, $m \in M_j$, each containing $0$, and a constant $C < \infty$ such that for every $j$,
every choice $a_m \in \mathcal{A}_{j,m}$ and all real coefficients $c = (c_m)_{m \in M_j}$,
\[
  \Bigl\| \sum_m c_m a_m \Bigr\| \leq C \, 2^{-j\delta} \|c\|_{\ell_p}
\]
and, for every $k \geq j$,
\[
  E_{2^k}\Bigl(\sum_m c_m a_m\Bigr) \leq 2^{-(k-j)\sigma} \, C \, 2^{-j\delta} \|c\|_{\ell_p} .
\]
For coefficients $\lambda = (\lambda_{j,m})$ and $s, q$ as before,
$\|\lambda\|_{b^s_{p,q}} = \|(2^{j(s-\delta)} \|\lambda_{j,\cdot}\|_{\ell_p})_j\|_{\ell_q}$.
\end{definition}

On a domain of dimension $d$ in $L_p$, the positions are the dyadic cubes of side $2^{-j}$ and
$\delta = d/p$. The size bound expresses the normalisation and the bounded overlap of the
supports, the decay bound the smoothness $\sigma$ of the atoms relative to the scheme, and
$b^s_{p,q}$ is the sequence space of Schneider and Vyb\'{\i}ral~\cite{SchneiderVybiral}.

\begin{theorem}[Atomic synthesis]
\label{thm:atomic-synthesis}
\lean{ApproximationScheme.AtomicSystem.approxNormDyadic_le_bNorm,
  ApproximationScheme.AtomicSystem.memApprox_of_bNorm_lt_top,
  ApproximationScheme.AtomicSystem.synthConst, ApproximationScheme.AtomicSystem.synthConst_ne_top}
\leanok
\uses{def:atomic-system, thm:synthesis-decay}
Let $A$ be a linear scheme with an atomic system as in \Cref{def:atomic-system}, let
$0 < s < \sigma$ and $1 \leq q \leq \infty$. Let $a_{j,m} \in \mathcal{A}_{j,m}$, let
$\lambda_{j,m} \in \R$, and let $f = \sum_j \sum_m \lambda_{j,m} a_{j,m}$, the series over $j$
converging in $X$. Then
\[
  \|f\|^{\mathrm{dy}}_{A^s_q} \leq \bigl(2 (1 - 2^{-s})^{-1} + (1 - 2^{-(\sigma - s)})^{-1}\bigr)
  \, C \, \|\lambda\|_{b^s_{p,q}} .
\]
\end{theorem}

\begin{proof}
\leanok
Apply \Cref{thm:synthesis-decay} to the levels $g_j = \sum_m \lambda_{j,m} a_{j,m}$ with
$\varepsilon_j = C \, 2^{-j\delta} \|\lambda_{j,\cdot}\|_{\ell_p}$; then
$2^{js} \varepsilon_j = C \, 2^{j(s - \delta)} \|\lambda_{j,\cdot}\|_{\ell_p}$.
\end{proof}

\begin{lemma}[Block decompositions with bounds]
\label{lem:exists-blocks}
\lean{ApproximationScheme.exists_blocks, ApproximationScheme.DyadicApprox.hasSum_blocks,
  ApproximationScheme.DyadicApprox.lqNorm_blocks_le,
  ApproximationScheme.DyadicApprox.tsum_enorm_blocks_ne_top,
  ApproximationScheme.DyadicApprox.blocks_mem, ApproximationScheme.DyadicApprox.sum_range_blocks,
  ApproximationScheme.DyadicApprox.blockDecompNorm_le}
\leanok
\uses{thm:block-decomposition}
Let $X$ be a Banach space, $s > 0$, $1 \leq q \leq \infty$ and $f \in A^s_q$. There are
$g_k \in A_{2^k}$ with $f = \sum_k g_k$, $\sum_k \|g_k\| < \infty$ and
$\|(2^{ks} \|g_k\|)_k\|_{\ell_q} \leq (3 \cdot 2^s + 4 \cdot 4^s) \|f\|^{\mathrm{dy}}_{A^s_q}$.
\end{lemma}

\begin{proof}
\leanok
These are the blocks of near-best approximants from the proof of
\Cref{thm:block-decomposition}.
\end{proof}

\begin{theorem}[Atomic analysis]
\label{thm:atomic-analysis}
\lean{ApproximationScheme.AtomicSystem.StableExpansion,
  ApproximationScheme.AtomicSystem.exists_atomic_decomposition,
  ApproximationScheme.AtomicSystem.memApprox_iff_exists_atomic}
\leanok
\uses{def:atomic-system, lem:exists-blocks, thm:atomic-synthesis}
Let $X$ be a Banach space with an atomic system for $A$, and let $C' \geq 0$ be such that every
$x \in A_{2^j}$ is $\sum_m c_m a_m$ with $a_m \in \mathcal{A}_{j,m}$ and
$\|c\|_{\ell_p} \leq C' \, 2^{j\delta} \|x\|$. Let $s > 0$ and $1 \leq q \leq \infty$. Then
every $f \in A^s_q$ is $\sum_j \sum_m \lambda_{j,m} a_{j,m}$ with $a_{j,m} \in
\mathcal{A}_{j,m}$ and
\[
  \|\lambda\|_{b^s_{p,q}} \leq C' (3 \cdot 2^s + 4 \cdot 4^s) \|f\|^{\mathrm{dy}}_{A^s_q} .
\]
If moreover $A$ is linear, $C' < \infty$ and $s < \sigma$, then $f \in A^s_q$ if and only if
$f$ is such an atomic sum with $\|\lambda\|_{b^s_{p,q}} < \infty$.
\end{theorem}

\begin{proof}
\leanok
Expand every block of \Cref{lem:exists-blocks} in atoms of its level; then
$2^{j(s-\delta)} \|\lambda_{j,\cdot}\|_{\ell_p} \leq C' \, 2^{js} \|g_j\|$. The converse is
\Cref{thm:atomic-synthesis}.
\end{proof}

Local bases on dyadic grids, such as the Haar functions, have the expansion property of
\Cref{thm:atomic-analysis}.

\chapter{Traces and intrinsic norms}
\label{ch:traces-abstract}

Throughout, $A$ is an approximation scheme in $X$ and $B$ one in $Y$. Two questions about
function spaces on domains become statements about approximation spaces: the restriction of
functions to a hypersurface, a map defined only on smooth functions, and the comparison of the
intrinsic norm of a domain with the norm of restrictions.

\begin{definition}[Trace bound]
\label{def:trace-bound}
\lean{ApproximationScheme.TraceBound, ApproximationScheme.TraceBound.ofSubmodule,
  ApproximationScheme.TraceBound.ofSubmodule_T, ApproximationScheme.traceConst,
  ApproximationScheme.traceConst_ne_top, ApproximationScheme.TraceBound.enorm_apply_le,
  ApproximationScheme.TraceBound.enorm_apply_sub_succ_le,
  ApproximationScheme.TraceBound.edist_apply_succ_le}
\leanok
\uses{def:scheme}
Let $\beta \in \R$. A \emph{trace bound} of exponent $\beta$ is a linear map $T : X \to Y$ with
$T(A_n) \subseteq B_n$ for every $n$ and a constant $C < \infty$ with
$\|T a\| \leq C n^\beta \|a\|$ for all $n \geq 1$ and $a \in A_n$.
\end{definition}

Only the values of $T$ on $\bigcup_n A_n$ matter, and a linear map defined on the span of the
approximants extends linearly to $X$. The model is the restriction of functions to a
hyperplane, defined on continuous functions and unbounded on $L_p$, with $\beta = 1/p$.

\begin{lemma}[The trace]
\label{lem:trace-ext}
\lean{ApproximationScheme.TraceBound.ext, ApproximationScheme.canonicalApprox,
  ApproximationScheme.TraceBound.tendsto_ext, ApproximationScheme.TraceBound.tendsto_edist_apply,
  ApproximationScheme.TraceBound.cauchySeq_apply, ApproximationScheme.TraceBound.ext_eq_of_mem,
  ApproximationScheme.TraceBound.ext_add, ApproximationScheme.TraceBound.ext_smul,
  ApproximationScheme.TraceBound.extLinear, ApproximationScheme.TraceBound.enorm_ext_le,
  ApproximationScheme.mem_iUnionSubmodule,
  ApproximationScheme.DyadicApprox.tendsto_dyadicWeight_mul_enorm_sub,
  ApproximationScheme.tsum_canonical_ne_top, ApproximationScheme.TraceBound.tendsto_canonical,
  ApproximationScheme.TraceBound.ext_zero}
\leanok
\uses{def:trace-bound, lem:near-best, def:dyadic-norm}
Let $T$ be a trace bound of exponent $\beta \geq 0$ and let $Y$ be a Banach space. For every
$f \in A^\beta_1$ there is $\operatorname{Tr} f \in Y$ such that $T c_k \to \operatorname{Tr} f$
for all $\gamma \in \N$ and all $c_k \in A_{2^{k+\gamma}}$ with
$2^{k\beta} \|f - c_k\| \to 0$. The map $\operatorname{Tr}$ is linear on $A^\beta_1$, agrees
with $T$ on every $A_m$, and
$\|\operatorname{Tr} f\| \leq (3C + 6C \cdot 4^\beta) \|f\|^{\mathrm{dy}}_{A^\beta_1}$.
\end{lemma}

\begin{proof}
\leanok
For near-best approximants $a_k$, $\|T a_{k+1} - T a_k\| \leq C \, 2^{(k+2)\beta} \cdot
2 (E_{2^k}(f) + t_k)$, which is summable; so $T a_k$ is a Cauchy sequence. Two admissible
sequences differ by $c_k - c'_k \in A_{2^{k+\gamma+\gamma'+1}}$ with
$\|T(c_k - c'_k)\| \lesssim 2^{k\beta} (\|f - c_k\| + \|f - c'_k\|) \to 0$. Linearity follows
by applying this to $a_k(f) + a_k(g)$ and $r a_k(f)$; on $A_m$ take $c_k = a$. The bound is the
telescoping sum, as in \Cref{thm:lift}.
\end{proof}

\begin{theorem}[Trace transfer]
\label{thm:trace-transfer}
\lean{ApproximationScheme.TraceBound.approxNormDyadic_ext_le,
  ApproximationScheme.TraceBound.memApprox_ext, ApproximationScheme.traceTransferConst,
  ApproximationScheme.TraceBound.edist_apply_ext_le, ApproximationScheme.traceTransferConst_ne_top,
  ApproximationScheme.TraceBound.enorm_ext_le_of_dyadicApprox}
\leanok
\uses{lem:trace-ext, lem:hardy}
Let $T$ be a trace bound of exponent $\beta$ with constant $C$, let $Y$ be a Banach space, and
let $0 \leq \beta < s$ and $1 \leq q \leq \infty$. For every $f \in A^s_q$,
\[
  \|\operatorname{Tr} f\|^{\mathrm{dy}}_{B^{s-\beta}_q} \leq
  \bigl(3C + 6C \cdot 4^\beta (1 - 2^{-(s-\beta)})^{-1} + 4C \cdot 4^\beta\bigr)
  \|f\|^{\mathrm{dy}}_{A^s_q} .
\]
\end{theorem}

\begin{proof}
\leanok
Since $T a_k \in B_{2^k}$, $E^B_{2^k}(\operatorname{Tr} f) \leq \sum_{i \geq 0}
\|T a_{k+i+1} - T a_{k+i}\|$. The proof of \Cref{thm:nikolskii-transfer} then applies word for
word; that theorem is the case $T = \iota^{-1}$ on the approximants.
\end{proof}

\begin{theorem}[Surjectivity of the trace]
\label{thm:trace-surjective}
\lean{ApproximationScheme.TraceBound.RightInverseBound, ApproximationScheme.shiftSeq,
  ApproximationScheme.shiftSeq_add, ApproximationScheme.shiftSeq_of_lt,
  ApproximationScheme.hasSum_shiftSeq, ApproximationScheme.shiftSeq_mem,
  ApproximationScheme.lqNorm_shiftSeq_le, ApproximationScheme.TraceBound.ext_eq_of_hasSum_lift,
  ApproximationScheme.TraceBound.exists_ext_eq, ApproximationScheme.traceSurjConst,
  ApproximationScheme.traceSurjConst_ne_top}
\leanok
\uses{thm:trace-transfer, lem:exists-blocks, thm:block-decomposition}
Let $X$ and $Y$ be Banach spaces, $T$ a trace bound of exponent $\beta \geq 0$, and let
$\gamma \in \N$ and $C' < \infty$ be such that every $b \in B_{2^k}$ is $T a$ for some
$a \in A_{2^{k+\gamma}}$ with $\|a\| \leq C' \, 2^{-k\beta} \|b\|$. Let $s > 0$ and
$1 \leq q \leq \infty$. For every $g \in B^s_q$ there is $f \in A^{s+\beta}_q$ with
$\operatorname{Tr} f = g$ and
\[
  \|f\|^{\mathrm{dy}}_{A^{s+\beta}_q} \leq 2 (1 - 2^{-(s+\beta)})^{-1} \, 2^{\gamma(s+\beta)}
  \, C' \, (3 \cdot 2^s + 4 \cdot 4^s) \, \|g\|^{\mathrm{dy}}_{B^s_q} .
\]
\end{theorem}

\begin{proof}
\leanok
Lift the blocks $b_j \in B_{2^j}$ of \Cref{lem:exists-blocks} for $g$ to $a_j \in
A_{2^{j+\gamma}}$ and put $f = \sum_j a_j$. The blocks $a_j$, placed at level $j + \gamma$,
form a block decomposition of $f$, and $2^{(j+\gamma)(s+\beta)} \|a_j\| \leq
2^{\gamma(s+\beta)} C' \, 2^{js} \|b_j\|$, so \Cref{thm:block-decomposition} gives the bound.
The partial sums $c_K = \sum_{j < K} a_j$ satisfy $2^{K\beta} \|f - c_K\| \leq C' \sum_{m \geq 0}
\|b_{m+K}\| \to 0$, so $T c_K \to \operatorname{Tr} f$ by \Cref{lem:trace-ext}, while
$T c_K = \sum_{j < K} b_j \to g$.
\end{proof}

This is where block decompositions are used: they turn the right inverse on approximants into
a right inverse of the trace.

\begin{definition}[Approximation-preserving extension]
\label{def:approx-extension}
\lean{ApproximationScheme.ApproxExtension, ApproximationScheme.restrictionNorm}
\leanok
\uses{def:error}
Let $R : X \to Y$ be a bounded linear map. An \emph{approximation-preserving extension} is a
map $E : Y \to X$, not necessarily linear, with $R(E g) = g$, $\|E g\| \leq C \|g\|$ and
$E^A_{2^{k+\gamma}}(E g) \leq C \, E^B_{2^k}(g)$ for all $g \in Y$ and $k \in \N$, for some
$C < \infty$ and $\gamma \in \N$. The \emph{restriction quasi-norm} of $g \in Y$ is
$\inf\{\|f\|^{\mathrm{dy}}_{A^s_q} : R f = g\}$.
\end{definition}

\begin{theorem}[Intrinsic and restriction quasi-norms]
\label{thm:retract}
\lean{ApproximationScheme.approxNormDyadic_extension_le, ApproximationScheme.extensionConst,
  ApproximationScheme.approxNormDyadic_le_iInf_restrict,
  ApproximationScheme.restrictionNorm_le_approxNormDyadic,
  ApproximationScheme.memApprox_iff_exists_restrict, ApproximationScheme.extensionConst_ne_top}
\leanok
\uses{def:approx-extension, lem:approx-functorial}
Let $R : X \to Y$ be a bounded linear map with $R(A_n) \subseteq B_n$ for every $n$, with an
approximation-preserving extension $E$. Let $s \geq 0$ and $1 \leq q \leq \infty$. Then for
every $g \in Y$,
\[
  \|g\|^{\mathrm{dy}}_{B^s_q} \leq \inf_{R f = g} \|R\| \, \|f\|^{\mathrm{dy}}_{A^s_q}
  \qquad \text{and} \qquad
  \|E g\|^{\mathrm{dy}}_{A^s_q} \leq C \bigl(1 + (\gamma + 1) 2^{\gamma s}\bigr)
  \|g\|^{\mathrm{dy}}_{B^s_q} ,
\]
and $g \in B^s_q$ if and only if $g = R f$ for some $f \in A^s_q$.
\end{theorem}

\begin{proof}
\leanok
The first inequality is \Cref{lem:approx-functorial}. For the second, the terms $k < \gamma$
are at most $2^{\gamma s} C \|g\|$, and for $k = j + \gamma$,
$2^{ks} E^A_{2^k}(E g) \leq 2^{\gamma s} C \, 2^{js} E^B_{2^j}(g)$.
\end{proof}

For a domain $\Omega \subseteq \R^d$, $X = L_p(\R^d)$, $Y = L_p(\Omega)$ and $R$ the restriction,
the quasi-norm of $B^s_q$ uses only $g$ on $\Omega$: it is intrinsic. The theorem reduces the
existence of an intrinsic norm to an approximation-preserving extension, which DeVore and
Sharpley~\cite{DeVoreSharpley} construct on Lipschitz domains from near-best local
polynomials. The extension may be
nonlinear on purpose: in parts of the quasi-Banach range no linear extension exists.

\chapter{Moduli of smoothness and approximation}
\label{ch:moduli-abstract}

Classical function spaces are defined by moduli of smoothness, approximation spaces by errors
of approximation. A Jackson inequality and an inverse inequality make them coincide. This
chapter states this once for every scheme; the examples are the trigonometric polynomials, the
piecewise polynomials, splines and wavelets.

\begin{definition}[Modulus compatible with a scheme]
\label{def:modulus-data}
\lean{ApproximationScheme.ModulusData, ApproximationScheme.ModulusData.modNorm,
  ApproximationScheme.ModulusData.modulus_zero, ApproximationScheme.ModulusData.modulus_sum_le}
\leanok
\uses{def:scheme, def:error}
Let $\theta \in \R$. A \emph{modulus compatible with $A$ of inverse order $\theta$} is a family
$\omega(f, k) \in [0, \infty]$, $f \in X$, $k \in \N$ (the modulus of smoothness at the scale
$2^{-k}$), together with a level $k_0 \in \N$ and constants $C_0, C_J, C_B < \infty$ such that for
all $f, g \in X$ and $j, k \in \N$
\[
  \omega(f + g, k) \leq \omega(f, k) + \omega(g, k), \qquad \omega(f, k) \leq C_0 \|f\|,
\]
\[
  E_{2^k}(f) \leq C_J \, \omega(f, k) \quad (k \geq k_0), \qquad
  \omega(a, k) \leq C_B \, 2^{-(k - j)_+ \theta} \|a\| \quad (a \in A_{2^j}) .
\]
The \emph{classical quasi-norm} is
$\|f\|^{\omega}_{s,q} = \|f\| + \|(2^{ks} \omega(f, k))_k\|_{\ell_q}$.
\end{definition}

\begin{theorem}[Coincidence of approximation and moduli]
\label{thm:modulus-coincidence}
\lean{ApproximationScheme.ModulusData.directConst,
  ApproximationScheme.ModulusData.directConst_ne_top,
  ApproximationScheme.ModulusData.approxNormDyadic_le_modNorm,
  ApproximationScheme.ModulusData.dyadicWeight_neg_sub_two_le,
  ApproximationScheme.ModulusData.inverseConst, ApproximationScheme.ModulusData.inverseConst_ne_top,
  ApproximationScheme.ModulusData.modulus_le_of_dyadicApprox,
  ApproximationScheme.ModulusData.dyadicWeight_mul_modulus_le_of_dyadicApprox,
  ApproximationScheme.ModulusData.modNorm_le_approxNormDyadic,
  ApproximationScheme.ModulusData.memApprox_iff_modNorm_lt_top}
\leanok
\uses{def:modulus-data, def:dyadic-norm, lem:hardy}
Let $\omega$ be a modulus compatible with $A$ of inverse order $\theta$ and $1 \leq q \leq
\infty$. Then for every $f \in X$ and $s \geq 0$,
$\|f\|^{\mathrm{dy}}_{A^s_q} \leq (1 + \sum_{k < k_0} 2^{ks} + C_J) \|f\|^{\omega}_{s,q}$, and
for $0 \leq s < \theta$,
\[
  \|f\|^{\omega}_{s,q} \leq \bigl(1 + 3 C_0 + C_B (3 + 6 \cdot 4^{\theta})
  (1 - 2^{-(\theta - s)})^{-1}\bigr) \|f\|^{\mathrm{dy}}_{A^s_q} .
\]
In particular $f \in A^s_q$ if and only if $\|f\|^{\omega}_{s,q} < \infty$, for $0 \leq s < \theta$.
\end{theorem}

\begin{proof}
\leanok
The first bound is Jackson's inequality for $k \geq k_0$ and $E_{2^k}(f) \leq \|f\|$ for
$k < k_0$. For the second,
let $a_k \in A_{2^k}$ be near-best approximants and write
$f = (f - a_k) + a_0 + \sum_{j < k} (a_{j+1} - a_j)$, where $a_{j+1} - a_j \in A_{2^{j+2}}$. Then
\begin{multline*}
  2^{ks} \omega(f, k) \leq C_0 \, 2^{ks} \|f - a_k\| + C_B \, 2^{-k(\theta - s)} \|a_0\| \\
  + 2 C_B 4^{\theta} \, 2^{-k(\theta - s)} \sum_{j \leq k} 2^{j(\theta - s)} \,
  2^{js} (E_{2^j}(f) + t_j),
\end{multline*}
and the last sum is bounded in $\ell_q$ by Hardy's inequality (\Cref{lem:hardy}) with exponent
$\theta - s > 0$.
\end{proof}

\chapter{Besov spaces and their embeddings}
\label{ch:besov}

Throughout, $\Omega$ is a topological space with its Borel $\sigma$-algebra and $\mu$ is a
finite measure on it. For $1 \leq p \leq \infty$, $L_p = L_p(\mu)$ is the space of
complex-valued $p$-integrable functions, and $C_b(\Omega)$ is the space of bounded continuous
complex-valued functions with the supremum norm.

\section{Three facts on \texorpdfstring{$L_p$}{L\_p}}

\begin{lemma}[The inclusion $L_r \hookrightarrow L_p$]
\label{lem:inclusion}
\lean{MeasureTheory.Lp.inclusionOfExponentLe, MeasureTheory.Lp.inclusionConst,
  MeasureTheory.Lp.coeFn_inclusionOfExponentLe, MeasureTheory.Lp.norm_inclusionOfExponentLe_le,
  MeasureTheory.Lp.norm_inclusionOfExponentLe_apply_le_of_prob,
  MeasureTheory.Lp.inclusionOfExponentLe_injective, MeasureTheory.Lp.inclusionOfExponentLe_comp,
  MeasureTheory.Lp.inclusionOfExponentLeLinear, MeasureTheory.Lp.coeFn_inclusionOfExponentLeLinear,
  MeasureTheory.Lp.inclusionConst_nonneg, MeasureTheory.Lp.inclusionConst_of_isProbabilityMeasure,
  MeasureTheory.Lp.norm_inclusionOfExponentLeLinear_apply_le,
  MeasureTheory.Lp.norm_inclusionOfExponentLe_apply_le,
  MeasureTheory.Lp.inclusionOfExponentLe_self}
\leanok
Let $1 \leq p \leq r \leq \infty$. The identity of functions defines an injective bounded
linear map $\iota_{r,p} : L_r \to L_p$ with $\|\iota_{r,p}\| \leq
\mu(\Omega)^{1/p - 1/r}$, and $\iota_{t,p} \circ \iota_{r,t} = \iota_{r,p}$ for every
$p \leq t \leq r$.
\end{lemma}

\begin{proof}
\leanok
H\"older's inequality gives $\|f\|_p \leq \mu(\Omega)^{1/p - 1/r} \|f\|_r$. A function that
vanishes almost everywhere does so in both spaces.
\end{proof}

\begin{lemma}[Interpolation between $L_p$ and $L_\infty$]
\label{lem:lp-interpolation}
\lean{ENNReal.one_sub_toReal_div_nonneg, MeasureTheory.eLpNorm_le_rpow_mul_eLpNorm_rpow,
  MeasureTheory.toReal_eLpNorm_le_rpow_mul, MeasureTheory.MemLp.mono_exponent_of_enorm_le}
\leanok
Let $0 < p \leq r < \infty$ and $M \geq 0$. For every measurable $f$ with $|f| \leq M$
almost everywhere,
\[
  \|f\|_r \leq M^{1 - p/r} \, \|f\|_p^{p/r} .
\]
For every $0 < p \leq r \leq \infty$ and every $f \in L_p$ with $|f| \leq M < \infty$ almost
everywhere, $f \in L_r$.
\end{lemma}

\begin{proof}
\leanok
Integrate $|f|^r \leq M^{r-p} |f|^p$. The exponent $1 - p/r$ is nonnegative, so the right-hand
side is finite for $f \in L_p$; for $r = \infty$, $\|f\|_\infty \leq M$.
\end{proof}

\begin{lemma}[Bounded continuous functions in $L_\infty$]
\label{lem:bc-linfty}
\lean{BoundedContinuousFunction.norm_toLp_top, BoundedContinuousFunction.isometry_toLp_top,
  BoundedContinuousFunction.isClosed_range_toLp_top}
\leanok
Assume that $\mu(U) > 0$ for every nonempty open $U \subseteq \Omega$. Then
$\|a\|_{L_\infty} = \sup_{x \in \Omega} |a(x)|$ for every $a \in C_b(\Omega)$, and the classes
of bounded continuous functions form a closed subspace of $L_\infty$.
\end{lemma}

\begin{proof}
\leanok
If $|a(x)| > c$ for some $x$, the open set $\{|a| > c\}$ has positive measure, so
$\|a\|_{L_\infty} \geq c$. The map $C_b(\Omega) \to L_\infty$ is therefore an isometry, and
its range is complete, hence closed.
\end{proof}

\section{Schemes of functions and Besov spaces}

\begin{definition}[Scheme of functions]
\label{def:function-scheme}
\lean{FunctionScheme, FunctionScheme.toLp, FunctionScheme.mem_toLp, FunctionScheme.toLp_mem_toLp,
  FunctionScheme.inclusion_mem_toLp}
\leanok
\uses{def:scheme, def:comap}
A \emph{scheme of functions} on $\Omega$ is an approximation scheme $S$ in $C_b(\Omega)$. For
$1 \leq p \leq \infty$, the classes in $L_p$ of the functions in $S_n$ form an approximation
scheme $S^{(p)}$ in $L_p$.
\end{definition}

Taking the approximants continuous and bounded makes them elements of every $L_p$ at once,
so that one scheme serves all exponents.

\begin{lemma}[Compatibility with the inclusions]
\label{lem:function-scheme-inclusion}
\lean{FunctionScheme.image_toLp_of_le, FunctionScheme.comap_toLp_of_le,
  BoundedContinuousFunction.inclusionOfExponentLe_toLp,
  BoundedContinuousFunction.norm_toLp_eq_toReal_eLpNorm, FunctionScheme.mem_toLp_of_inclusion_mem,
  FunctionScheme.exists_toLp_of_mem}
\leanok
\uses{def:function-scheme, lem:inclusion}
Let $S$ be a scheme of functions and $1 \leq p \leq r \leq \infty$. For every $n \in \N$,
$\iota_{r,p}(S^{(r)}_n) = S^{(p)}_n$ and $\iota_{r,p}^{-1}(S^{(p)}_n) = S^{(r)}_n$.
\end{lemma}

\begin{proof}
\leanok
The map $\iota_{r,p}$ sends the class of $a$ in $L_r$ to its class in $L_p$, and it is
injective.
\end{proof}

\begin{definition}[Besov space]
\label{def:besov}
\lean{FunctionScheme.besovNorm, FunctionScheme.MemBesov, FunctionScheme.besovSpace,
  FunctionScheme.besovSpaceSup}
\leanok
\uses{def:function-scheme, def:approx-space, prop:approx-linear}
Let $S$ be a scheme of functions, $s \geq 0$ and $1 \leq p, q \leq \infty$. The
\emph{Besov space} is $\Besov{s}{p}{q}(\Omega, S) = A^s_q(L_p, S^{(p)})$, the approximation
space of $L_p$ for the scheme $S^{(p)}$. Its errors $E_n(f)_p$ are the errors of best
approximation of \Cref{def:error} for the scheme $S^{(p)}$.
\end{definition}

On the torus $\T^d$ with $S_n$ the trigonometric polynomials of degree $< n$, these are the
classical periodic Besov spaces (\Cref{thm:moduli-coincidence}).

\begin{proposition}[Elementary embeddings]
\label{prop:besov-elementary}
\lean{FunctionScheme.MemBesov.mono, FunctionScheme.besovNorm_inclusion_le,
  FunctionScheme.MemBesov.inclusion, FunctionScheme.MemBesov.of_sup}
\leanok
\uses{def:besov, thm:approx-embedding, lem:approx-functorial,
  lem:function-scheme-inclusion}
Let $S$ be a scheme of functions and $1 \leq p, q, q' \leq \infty$.
\begin{enumerate}
  \item If $0 \leq t$, and $t < s$ or $t = s$ and $q \leq q'$, then
    $\Besov{s}{p}{q} \subseteq \Besov{t}{p}{q'}$.
  \item If $p \leq r$, then $\iota_{r,p} f \in \Besov{s}{p}{q}$ and
    $\|\iota_{r,p} f\|_{\Besov{s}{p}{q}} \leq \|\iota_{r,p}\| \, \|f\|_{\Besov{s}{r}{q}}$ for
    every $f \in \Besov{s}{r}{q}$.
  \item The class in $L_p$ of every $a \in A^s_q(C_b(\Omega), S)$ lies in
    $\Besov{s}{p}{q}$.
\end{enumerate}
\end{proposition}

\begin{proof}
\leanok
The first item is \Cref{thm:approx-embedding}. The others are \Cref{lem:approx-functorial}
for $T = \iota_{r,p}$, using \Cref{lem:function-scheme-inclusion}, and for the map
$C_b(\Omega) \to L_p$, of norm at most $\mu(\Omega)^{1/p}$.
\end{proof}

\section{The Nikolskii inequality and the Besov embeddings}

\begin{definition}[Nikolskii property]
\label{def:nikolskii-property}
\lean{FunctionScheme.HasNikolskii, nikolskiiExponent, nikolskiiExponent_nonneg,
  one_le_of_nikolskiiExponent_lt, nikolskiiExponent_top, nikolskiiExponent_self,
  FunctionScheme.toReal_eLpNorm_top_le_norm}
\leanok
\uses{def:function-scheme}
Let $S$ be a scheme of functions and $D \in \R$. The scheme has the \emph{Nikolskii property}
of dimension $D$ if for every $1 \leq p \leq \infty$ there is $C \geq 0$ with
\[
  \sup_{x \in \Omega} |a(x)| \leq C \, n^{D/p} \, \|a\|_{L_p}
  \qquad \text{for all } n \geq 1 \text{ and } a \in S_n .
\]
For $1 \leq p \leq r \leq \infty$ put $\alpha(D, p, r) = D(1/p - 1/r)$.
\end{definition}

The Nikolskii property enters the theorems below as a hypothesis on the scheme; the
trigonometric polynomials on the torus have it (\Cref{thm:torus-nikolskii}).

\begin{lemma}[Nikolskii inequality]
\label{lem:nikolskii-inequality}
\lean{FunctionScheme.HasNikolskii.eLpNorm_le, FunctionScheme.HasNikolskii.liftBound}
\leanok
\uses{def:nikolskii-property, lem:lp-interpolation, def:lift-bound,
  lem:function-scheme-inclusion}
Let $S$ have the Nikolskii property of dimension $D$, and let $1 \leq p \leq r \leq \infty$.
There is $C \geq 0$ with $\|a\|_{L_r} \leq C \, n^{\alpha(D,p,r)} \|a\|_{L_p}$ for all
$n \geq 1$ and $a \in S_n$. In particular $\iota_{r,p}$ satisfies the lifting bound of
exponent $\alpha(D, p, r)$ for $S^{(p)}$.
\end{lemma}

\begin{proof}
\leanok
For $r = \infty$, $\|a\|_{L_\infty} \leq \sup |a|$. For $r < \infty$, \Cref{lem:lp-interpolation}
with $M = \sup|a| \leq C n^{D/p} \|a\|_p$ gives $\|a\|_r \leq C^{1-p/r} n^{(D/p)(1-p/r)}
\|a\|_p$, and $(D/p)(1 - p/r) = \alpha(D, p, r)$. The class of $a$ in $L_r$ lifts the class
in $L_p$.
\end{proof}

\begin{theorem}[Besov embeddings]
\label{thm:besov-embedding}
\lean{FunctionScheme.exists_inclusion_eq_of_nikolskii,
  FunctionScheme.MemBesov.exists_inclusion_eq_of_nikolskii,
  FunctionScheme.MemBesov.exists_inclusion_eq_top,
  FunctionScheme.MemBesov.exists_boundedContinuous}
\leanok
\uses{lem:nikolskii-inequality, thm:nikolskii-transfer, lem:function-scheme-inclusion,
  lem:approx-closure, lem:bc-linfty}
Let $S$ have the Nikolskii property of dimension $D \geq 0$, let $1 \leq p \leq r \leq
\infty$ and $1 \leq q \leq \infty$, and let $s > \alpha = \alpha(D, p, r)$.
\begin{enumerate}
  \item There is $K < \infty$ such that for every $f \in \Besov{s}{p}{q}$ there is $g \in L_r$
    with $\iota_{r,p} g = f$ and $\|g\|^{\mathrm{dy}}_{\Besov{s-\alpha}{r}{q}} \leq K
    \|f\|^{\mathrm{dy}}_{\Besov{s}{p}{q}}$.
  \item If $\mu(U) > 0$ for every nonempty open $U \subseteq \Omega$ and $s > D/p$, then every
    $f \in \Besov{s}{p}{q}$ is the class of a bounded continuous function.
\end{enumerate}
\end{theorem}

\begin{proof}
\leanok
For the first item apply \Cref{thm:nikolskii-transfer} to $\iota = \iota_{r,p}$, which is
injective and satisfies the lifting bound of exponent $\alpha$ by
\Cref{lem:nikolskii-inequality}; the preimage scheme is $S^{(r)}$ by
\Cref{lem:function-scheme-inclusion}, and $L_r$ is complete. For the second take
$r = \infty$: then $g \in \Besov{s - D/p}{\infty}{q}$ with $s - D/p > 0$, so $g$ lies in the
closure of $\bigcup_n S^{(\infty)}_n$ (\Cref{lem:approx-closure}), which is contained in the
closed set of classes of bounded continuous functions (\Cref{lem:bc-linfty}).
\end{proof}

The limiting case $s = \alpha$ is out of reach of this method, since the Nikolskii transfer
theorem needs $\alpha < s$.

\chapter{Sobolev spaces and their embeddings}
\label{ch:sobolev}

The setting is that of \Cref{ch:besov}. A Sobolev space enters through its Jackson and
Bernstein inequalities, as a smoothness space for the scheme in $L_p$. The definition below is a
hypothesis whose content is these two inequalities; its instances are the spaces $H^s$ with
$p = 2$ on the torus (\Cref{def:torus-sobolev}) and for a continuous spectral basis
(\Cref{thm:continuous-spectral-basis-besov}).

\begin{definition}[Sobolev space for a scheme]
\label{def:sobolev-scheme}
\lean{FunctionScheme.SobolevSpace}
\leanok
\uses{def:smoothness, def:function-scheme}
Let $S$ be a scheme of functions and $1 \leq p \leq \infty$. A \emph{Sobolev space} of order
$s > 0$ for $S$ in $L_p$ is a smoothness space of order $s$ for $S^{(p)}$: a normed space $V$
with an injective bounded linear map $\iota : V \to L_p$, Jackson's inequality
$E_n(\iota v)_p \leq C n^{-s} \|v\|$ and Bernstein's inequality, every $a \in S^{(p)}_n$ being
$\iota v$ with $\|v\| \leq C n^s \|a\|_p$.
\end{definition}

\begin{theorem}[Sobolev embeddings]
\label{thm:sobolev-embedding}
\lean{FunctionScheme.SobolevSpace.exists_eq_inclusion,
  FunctionScheme.SobolevSpace.exists_eq_inclusion',
  FunctionScheme.SobolevSpace.exists_boundedContinuous}
\leanok
\uses{def:sobolev-scheme, thm:besov-embedding, prop:smoothness}
Let $S$ have the Nikolskii property of dimension $D \geq 0$, and let $1 \leq p \leq r \leq
\infty$. Let $(V, \iota)$ be a Sobolev space of order $s$ for $S$ in $L_p$.
\begin{enumerate}
  \item Let $(V', \iota')$ be a Sobolev space of order $t$ for $S$ in $L_r$ with $V'$ a Banach
    space, and let $t - D/r < s - D/p$. For every $v \in V$ there is $v' \in V'$ with
    $\iota_{r,p}(\iota' v') = \iota v$.
  \item If $\mu(U) > 0$ for every nonempty open $U \subseteq \Omega$ and $s > D/p$, then
    $\iota v$ is the class of a bounded continuous function for every $v \in V$.
\end{enumerate}
\end{theorem}

\begin{proof}
\leanok
The condition $t - D/r < s - D/p$ is $t < s - \alpha(D, p, r)$, and it implies $s > \alpha$
since $t > 0$. By \Cref{prop:smoothness}, $\iota v \in \Besov{s}{p}{\infty}$. By
\Cref{thm:besov-embedding}, $\iota v = \iota_{r,p} g$ with $g \in
\Besov{s-\alpha}{r}{\infty}$, and $g \in \iota'(V')$ by \Cref{prop:smoothness}, as
$s - \alpha > t$. The second item is \Cref{thm:besov-embedding}(2) with $q = \infty$.
\end{proof}

The limiting cases $t - D/r = s - D/p$ and $s = D/p$ are not reached: the passage from $V$
to $\Besov{s}{p}{\infty}$ and back from $\Besov{t}{r}{1}$ to $V'$ loses the fine index at
both ends.

\chapter{Spectral schemes}
\label{ch:spectral}

Throughout, $H$ is a complex Hilbert space with a Hilbert basis $(e_i)_{i \in I}$, and
$\hat f(i) = \langle e_i, f \rangle$ are the coefficients of $f \in H$.

\begin{definition}[Spectral scheme]
\label{def:spectral}
\lean{SpectralScheme, SpectralScheme.span, SpectralScheme.toScheme, SpectralScheme.weight,
  SpectralScheme.span_zero, SpectralScheme.span_mono, SpectralScheme.mem_toScheme,
  SpectralScheme.weight_pos, SpectralScheme.one_le_weight, SpectralScheme.norm_ofReal_inv_weight,
  SpectralScheme.norm_ofReal_weight, SpectralScheme.ofReal_weight_ne_zero,
  SpectralScheme.inv_weight_le, SpectralScheme.weight_le}
\leanok
\uses{def:comap}
A \emph{level function} is a map $\lambda : I \to \N$ such that $\{i : \lambda(i) < n\}$ is
finite for every $n$. The sets $A_n = \spn\{e_i : \lambda(i) < n\}$ form an approximation
scheme in $H$, the \emph{spectral scheme} of $(e_i, \lambda)$.
\end{definition}

On the torus, $e_k$ is the Fourier basis and $\lambda(k) = \max_j |k_j|$. On a compact
manifold, $e_i$ are the eigenfunctions of the Laplacian and $\lambda$ a function of the
eigenvalue; \Cref{sec:continuous-spectral-bases} assumes an orthonormal basis of continuous
functions with levels, the structure this case provides.

\begin{definition}[Sobolev embedding of $\ell_2$]
\label{def:sobolev-embed}
\lean{SpectralScheme.sobolevEmbed, SpectralScheme.repr_sobolevEmbed,
  SpectralScheme.sobolevEmbed_injective, SpectralScheme.diagInvWeight,
  SpectralScheme.diagInvWeight_apply}
\leanok
\uses{def:spectral}
For $s \geq 0$ let $J_s : \ell_2(I) \to H$ be the bounded injective linear map
\[
  J_s c = \sum_{i \in I} (1 + \lambda(i))^{-s} c_i \, e_i .
\]
\end{definition}

\begin{lemma}[Range of the Sobolev embedding]
\label{lem:sobolev-range}
\lean{SpectralScheme.mem_range_sobolevEmbed_iff, SpectralScheme.sobolevEmbed_weighted}
\leanok
\uses{def:sobolev-embed}
Let $s \geq 0$. An element $f \in H$ lies in $J_s(\ell_2(I))$ if and only if
$\sum_i (1 + \lambda(i))^{2s} |\hat f(i)|^2 < \infty$.
\end{lemma}

\begin{proof}
\leanok
The preimage is $c_i = (1 + \lambda(i))^s \hat f(i)$.
\end{proof}

\begin{proposition}[Jackson and Bernstein inequalities]
\label{prop:spectral-jackson-bernstein}
\lean{SpectralScheme.norm_sub_truncate_le, SpectralScheme.jacksonBound, SpectralScheme.liftBound,
  SpectralScheme.toSmoothnessSpace, SpectralScheme.truncate,
  SpectralScheme.repr_eq_zero_of_mem_span, SpectralScheme.truncate_mem_span,
  SpectralScheme.repr_truncate, SpectralScheme.repr_sub_truncate}
\leanok
\uses{def:sobolev-embed, def:jackson, def:lift-bound, def:smoothness}
Let $s \geq 0$. For every $n \geq 1$ and $c \in \ell_2(I)$, $E_n(J_s c) \leq n^{-s} \|c\|$;
and every $a \in A_n$ is $J_s c$ for some $c$ with $\|c\| \leq n^s \|a\|$. For $s > 0$,
$(\ell_2(I), J_s)$ is a smoothness space of order $s$ for the spectral scheme, with
constants $1$.
\end{proposition}

\begin{proof}
\leanok
The truncation $P_n f = \sum_{\lambda(i) < n} \hat f(i) e_i$ lies in $A_n$, and by
Parseval's identity $\|J_s c - P_n J_s c\|^2 = \sum_{\lambda(i) \geq n} (1 +
\lambda(i))^{-2s} |c_i|^2 \leq n^{-2s} \|c\|^2$. An element $a \in A_n$ has $\hat a(i) = 0$
for $\lambda(i) \geq n$, and $c_i = (1 + \lambda(i))^s \hat a(i)$ satisfies
$|c_i| \leq n^s |\hat a(i)|$.
\end{proof}

\begin{proposition}[Spectral projections]
\label{prop:spectral-projections}
\lean{SpectralScheme.projectionScheme, SpectralScheme.truncateCLM,
  SpectralScheme.closure_iUnion_toScheme, SpectralScheme.repr_block, SpectralScheme.enorm_block_sq,
  SpectralScheme.truncate_add, SpectralScheme.truncate_smul, SpectralScheme.norm_truncate_le,
  SpectralScheme.truncateCLM_apply, SpectralScheme.truncate_of_mem,
  SpectralScheme.projectionScheme_P, SpectralScheme.iUnion_span, SpectralScheme.enorm_sq_eq_tsum}
\leanok
\uses{def:spectral, def:projection-scheme, def:blocks}
The truncations $P_n$ form a projection scheme with bound $1$ for the spectral scheme, the
union $\bigcup_n A_n$ is dense in $H$, and for every $k \in \N$ and $f \in H$
\[
  \Delta_k f = \sum_{2^k \leq \lambda(i) < 2^{k+1}} \hat f(i) \, e_i ,
  \qquad
  \|\Delta_k f\|^2 = \sum_{2^k \leq \lambda(i) < 2^{k+1}} |\hat f(i)|^2 .
\]
\end{proposition}

\begin{proof}
\leanok
Truncation is the orthogonal projection onto $A_n$. The union contains the span of the
basis, which is dense. The formulas are Parseval's identity.
\end{proof}

\begin{theorem}[Coefficient characterisation]
\label{thm:spectral-coefficients}
\lean{SpectralScheme.memApprox_iff_blockNorm_lt_top, SpectralScheme.memApprox_two_iff,
  SpectralScheme.mem_l2_iff, SpectralScheme.tsum_ite_block_eq, SpectralScheme.enorm_weight_mul_sq}
\leanok
\uses{prop:spectral-projections, thm:blocks, lem:sobolev-range}
Let $s > 0$ and $1 \leq q \leq \infty$. An element $f \in H$ lies in $A^s_q$ if and only if
\[
  \Bigl(\sum_{\lambda(i) = 0} |\hat f(i)|^2\Bigr)^{1/2} + \Bigl\| \Bigl(2^{ks}
    \Bigl(\sum_{2^k \leq \lambda(i) < 2^{k+1}} |\hat f(i)|^2\Bigr)^{1/2}\Bigr)_k
    \Bigr\|_{\ell_q} < \infty .
\]
It lies in $A^s_2$ if and only if $\sum_i (1 + \lambda(i))^{2s} |\hat f(i)|^2 < \infty$; that
is, $A^s_2 = J_s(\ell_2(I))$.
\end{theorem}

\begin{proof}
\leanok
The first statement is \Cref{thm:blocks} with the formulas of
\Cref{prop:spectral-projections}, since $P_1 f = \sum_{\lambda(i) = 0} \hat f(i) e_i$. For
$q = 2$ the square of the block norm is comparable with $\sum_i \omega_i |\hat f(i)|^2$,
where $\omega_i = 2^{2ks}$ on the block $2^k \leq \lambda(i) < 2^{k+1}$ and $\omega_i = 1$
for $\lambda(i) = 0$; and $\omega_i \leq (1 + \lambda(i))^{2s} \leq 4^s \omega_i$.
\end{proof}


\section{Continuous spectral bases}
\label{sec:continuous-spectral-bases}

Let $M$ be a topological space with a finite Borel measure $\mu$. On a compact Riemannian
manifold the eigenfunctions of the Laplace--Beltrami operator form an orthonormal basis of
$L_2(M)$; they are continuous by elliptic regularity, and the eigenvalues $\lambda_i$ tend to
infinity. Mathlib has neither the operator nor these facts. The statements below therefore assume
an orthonormal basis of continuous functions with levels, and every one of them is proved from
this assumption; the torus satisfies it (\Cref{prop:fourier-eigenbasis}). The assumption
consists of the basis, its levels and its continuous representatives alone.

\begin{definition}[Continuous spectral basis]
\label{def:continuous-spectral-basis}
\lean{ContinuousSpectralBasis, ContinuousSpectralBasis.basisSpan,
  ContinuousSpectralBasis.basisScheme,
  ContinuousSpectralBasis.basisSpan_zero, ContinuousSpectralBasis.basisSpan_mono,
  ContinuousSpectralBasis.basisFun_mem_basisSpan, ContinuousSpectralBasis.mem_basisScheme,
  ContinuousSpectralBasis.toLp_basisScheme_eq}
\leanok
\uses{def:spectral, def:function-scheme}
A \emph{continuous spectral basis} of $L_2(\mu)$ is a spectral scheme $(e_i, \lambda)$ in
$L_2(\mu)$ together with bounded continuous functions $\tilde e_i$ whose classes are the $e_i$.
The spaces $\mathcal{E}_n = \spn\{\tilde e_i : \lambda(i) < n\}$ of bounded continuous functions
form a scheme of functions on $M$, and its classes in $L_2(\mu)$ form the spectral scheme of
$(e_i, \lambda)$.
\end{definition}

The structure enters the theorems below as a hypothesis; the Fourier basis of the torus is an
instance (\Cref{prop:fourier-eigenbasis}). For the Laplace--Beltrami operator, $\tilde e_i$ are
the eigenfunctions and $\lambda(i) = \lfloor \sqrt{\lambda_i} \rfloor$. The Besov spaces
$\Besov{s}{p}{q}(M)$ of the scheme $(\mathcal{E}_n)$ are defined for every $1 \leq p \leq \infty$
by \Cref{def:besov}.

\begin{theorem}[Besov and Sobolev spaces of a continuous spectral basis]
\label{thm:continuous-spectral-basis-besov}
\lean{ContinuousSpectralBasis.memBesov_two_two_iff,
  ContinuousSpectralBasis.memBesov_two_iff_blockNorm_lt_top,
  ContinuousSpectralBasis.sobolevSpace, ContinuousSpectralBasis.sobolevSpace_order,
  ContinuousSpectralBasis.sobolevSpace_incl, ContinuousSpectralBasis.mem_range_sobolev_iff,
  ContinuousSpectralBasis.sobolev_memBesov_top, ContinuousSpectralBasis.exists_sobolev_of_memBesov}
\leanok
\uses{def:continuous-spectral-basis, thm:spectral-coefficients, prop:spectral-jackson-bernstein,
  def:sobolev-scheme}
Let $s > 0$ and $1 \leq q \leq \infty$. For $f \in L_2(\mu)$:
\begin{enumerate}
  \item $f \in \Besov{s}{2}{2}(M)$ if and only if
    $\sum_i (1 + \lambda(i))^{2s} |\langle e_i, f \rangle|^2 < \infty$;
  \item $f \in \Besov{s}{2}{q}(M)$ if and only if the block norm of \Cref{thm:spectral-coefficients}
    is finite;
  \item the space $H^s(M)$ of the $f$ with
    $\sum_i (1 + \lambda(i))^{2s} |\langle e_i, f \rangle|^2 < \infty$ is a Sobolev space of order
    $s$ for the scheme in $L_2$, with constants $1$; $H^s(M) \subseteq \Besov{s}{2}{\infty}(M)$, and
    $\Besov{t}{2}{q}(M) \subseteq H^s(M)$ for $t > s$.
\end{enumerate}
\end{theorem}

\begin{proof}
\leanok
The classes in $L_2$ of the scheme form the spectral scheme
(\Cref{def:continuous-spectral-basis}), so \Cref{thm:spectral-coefficients} and
\Cref{prop:spectral-jackson-bernstein} apply.
\end{proof}

\begin{definition}[Local Weyl bound]
\label{def:local-weyl}
\lean{ContinuousSpectralBasis.LocalWeylBound, ContinuousSpectralBasis.levelSet,
  ContinuousSpectralBasis.mem_levelSet}
\leanok
\uses{def:continuous-spectral-basis}
A continuous spectral basis satisfies the \emph{local Weyl bound} of dimension $D$ with
constant $C \geq 0$ if
\[
  \sum_{\lambda(i) < n} |\tilde e_i(x)|^2 \leq C n^D
  \qquad \text{for every } x \in M \text{ and every } n \geq 1 .
\]
\end{definition}

For the Laplace--Beltrami operator of a compact Riemannian manifold of dimension $d$ this is
H\"ormander's local Weyl law, with $D = d$. The bound enters the theorems below as a
hypothesis; the Fourier basis of the torus satisfies it (\Cref{prop:fourier-eigenbasis}).

\begin{theorem}[Nikolskii inequality from the local Weyl bound]
\label{thm:local-weyl-nikolskii}
\lean{ContinuousSpectralBasis.norm_sq_le_of_localWeyl,
  ContinuousSpectralBasis.norm_le_sqrt_mul_of_localWeyl,
  ContinuousSpectralBasis.norm_rpow_le_of_localWeyl, ContinuousSpectralBasis.norm_le_of_localWeyl,
  ContinuousSpectralBasis.exists_sum_of_mem_basisSpan, ContinuousSpectralBasis.norm_toLp_sum_sq}
\leanok
\uses{def:local-weyl, lem:lp-interpolation}
Let the continuous spectral basis satisfy the local Weyl bound of dimension $D$ with constant
$C$. For every $n \geq 1$, every $a \in \mathcal{E}_n$ and every $1 \leq p \leq 2$,
\[
  \|a\|_\infty^2 \leq C n^D \|a\|_2^2 ,
  \qquad
  \|a\|_\infty^p \leq C n^D \|a\|_p^p .
\]
\end{theorem}

\begin{proof}
\leanok
Write $a = \sum_{\lambda(i) < n} c_i \tilde e_i$. By the Cauchy--Schwarz inequality and the local
Weyl bound, $|a(x)|^2 \leq \sum_i |c_i|^2 \sum_{\lambda(i) < n} |\tilde e_i(x)|^2 \leq C n^D
\sum_i |c_i|^2$, and $\sum_i |c_i|^2 = \|a\|_2^2$ by Parseval's identity. For $p \leq 2$,
$\|a\|_2^2 \leq \|a\|_\infty^{2-p} \|a\|_p^p$ (\Cref{lem:lp-interpolation}); dividing by
$\|a\|_\infty^{2-p}$ gives the second inequality.
\end{proof}

For $p > 2$ the inequality $\|a\|_\infty \leq C' n^{D/p} \|a\|_p$ does not follow in this way:
the proof on the torus raises polynomials to powers (\Cref{thm:torus-nikolskii}), which lie in
$\mathcal{E}_{mn}$ (\Cref{lem:trig-products}); for a general basis, products of functions of
$\mathcal{E}_n$ need not lie in any $\mathcal{E}_m$. The embeddings of \Cref{thm:besov-embedding}
hold for the scheme $(\mathcal{E}_n)$ whenever it has the Nikolskii property of
\Cref{def:nikolskii-property}.
The embedding into continuous functions needs only the local Weyl bound.

\begin{theorem}[Continuity of $H^s(M)$]
\label{thm:local-weyl-continuity}
\lean{ContinuousSpectralBasis.exists_continuous_of_localWeyl, ContinuousSpectralBasis.partialSum,
  ContinuousSpectralBasis.partialSum_mem_basisSpan, ContinuousSpectralBasis.toLp_partialSum,
  ContinuousSpectralBasis.norm_truncate_sub_truncate_le,
  ContinuousSpectralBasis.norm_partialSum_succ_sub_le, Real.natCast_two_pow_rpow}
\leanok
\uses{thm:local-weyl-nikolskii, thm:continuous-spectral-basis-besov}
Let the continuous spectral basis satisfy the local Weyl bound of dimension $D$, and let
$s > 0$ with $2s > D$. Every $f \in L_2(\mu)$ with
$\sum_i (1 + \lambda(i))^{2s} |\langle e_i, f \rangle|^2 < \infty$ is the class of a bounded
continuous function.
\end{theorem}

\begin{proof}
\leanok
Let $f = J_s c$ and $S_n f = \sum_{\lambda(i) < n} \langle e_i, f \rangle \tilde e_i$, a bounded
continuous function whose class is the truncation $P_n f$. The difference
$S_{2^{k+1}} f - S_{2^k} f$ lies in $\mathcal{E}_{2^{k+1}}$, and its $L_2$ norm is at most
$\|f - P_{2^k} f\|_2 \leq 2^{-ks} \|c\|$ (\Cref{prop:spectral-jackson-bernstein}). By
\Cref{thm:local-weyl-nikolskii}
\[
  \|S_{2^{k+1}} f - S_{2^k} f\|_\infty \leq C^{1/2} 2^{(k+1)D/2} 2^{-ks} \|c\|
  = C^{1/2} 2^{D/2} \|c\| \, (2^{D/2 - s})^k ,
\]
a geometric sequence since $D/2 < s$. So $(S_{2^k} f)_k$ converges uniformly to a bounded
continuous $g$, and the class of $g$ is the $L_2$ limit $\lim_k P_{2^k} f = f$.
\end{proof}
